\chapter{Historical audit and superseded adapter experiments}
The archived audit is part of the work completed, but its historical 80-image results and validation-only temporal ablations are not the corrected test protocol. The fifteen findings explain why later source labels, metrics, cached identities, checkpoints and split rules were strengthened. Their titles are preserved as audit observations; the accompanying source contains the reproductions and detailed limitations.
\begingroup\footnotesize
\setlength{\LTpre}{5pt}\setlength{\LTpost}{8pt}
\begin{longtable}{@{}>{\raggedright\arraybackslash}p{0.0737\linewidth}>{\raggedright\arraybackslash}p{0.1032\linewidth}>{\raggedright\arraybackslash}p{0.6632\linewidth}@{}}
\caption[Historical audit]{Fifteen archived audit findings; original reproductions are in the registered audit source.}\label{tab:historical_audit}\\
\toprule
\textbf{ID} & \textbf{Severity} & \textbf{Historical issue} \\
\midrule\endfirsthead
\multicolumn{3}{l}{\textit{Table \thetable\ continued}}\\
\toprule
\textbf{ID} & \textbf{Severity} & \textbf{Historical issue} \\
\midrule\endhead
\midrule\multicolumn{3}{r}{\textit{Continued on next page}}\\\endfoot
\bottomrule\endlastfoot
F01 & High & fractional positives are silently converted to negatives \\
F02 & High & reported optimal thresholds do not match their metrics \\
F03 & High & the CPU best-weight snapshot changes during further training \\
F04 & High & validation is reused for model/epoch/branch/threshold selection \\
F05 & High & dataset labels do not fully match the specified events \\
F06 & High & source and negative distribution confound the image comparison \\
F07 & High & the requested five-anomaly coverage is incomplete \\
F08 & Medium & abstentions are excluded from most headline metrics \\
F09 & Medium & macro-F1 has a changing class denominator \\
F10 & Medium & bootstrap ignores the group structure and does not establish a winner \\
F11 & Medium & padded frames influence real predictions \\
F12 & Medium & session labels/order and feature-cache reuse need stronger provenance \\
F13 & Medium & notebook reuse can silently show stale results \\
F14 & Medium & the old VadCLIP wrapper is not a working upstream integration \\
F15 & Medium & hardware/runtime tests do not select the best deployment pipeline \\
\end{longtable}\endgroup


The following older six-run VadCLIP suite varied graph use, attention window and representative length. It used a 12-session train/four-session validation split with no separate final test. The archived audit identified fractional-label truncation, threshold inconsistency and CPU snapshot/padding issues. Representative target counts differ with visual length. Therefore these archived AP/F1 entries document what was run, not valid evidence that one architecture generalizes better. They remain separate from the corrected 1,056-frame held-out test in Chapter~\ref{ch:baseline}.
\begingroup\footnotesize
\setlength{\LTpre}{5pt}\setlength{\LTpost}{8pt}
\begin{longtable}{@{}>{\raggedright\arraybackslash}p{0.1583\linewidth}>{\raggedright\arraybackslash}p{0.0852\linewidth}>{\raggedright\arraybackslash}p{0.0974\linewidth}>{\raggedright\arraybackslash}p{0.0852\linewidth}>{\raggedright\arraybackslash}p{0.0852\linewidth}>{\raggedright\arraybackslash}p{0.0852\linewidth}>{\raggedright\arraybackslash}p{0.1096\linewidth}>{\raggedright\arraybackslash}p{0.1339\linewidth}@{}}
\caption[Historical vadclip ablation]{Superseded validation-only adapter records; target/compression and implementation defects preclude a clean architecture claim.}\label{tab:historical_vadclip_ablation}\\
\toprule
\textbf{Run} & \textbf{Length} & \textbf{Window} & \textbf{Graph} & \textbf{Reps} & \textbf{+ reps} & \textbf{Old AP} & \textbf{Old opt F1} \\
\midrule\endfirsthead
\multicolumn{8}{l}{\textit{Table \thetable\ continued}}\\
\toprule
\textbf{Run} & \textbf{Length} & \textbf{Window} & \textbf{Graph} & \textbf{Reps} & \textbf{+ reps} & \textbf{Old AP} & \textbf{Old opt F1} \\
\midrule\endhead
\midrule\multicolumn{8}{r}{\textit{Continued on next page}}\\\endfoot
\bottomrule\endlastfoot
full & 256 & 8 & on & 383 & 115 & 0.5253 & 0.4891 \\
no graph & 256 & 8 & off & 383 & 115 & 0.5105 & 0.4934 \\
win 1 & 256 & 1 & on & 383 & 115 & 0.5419 & 0.4635 \\
win 16 & 256 & 16 & on & 383 & 115 & 0.5361 & 0.4870 \\
vl 128 & 128 & 8 & on & 255 & 60 & 0.6374 & 0.5778 \\
vl 512 & 512 & 8 & on & 639 & 276 & 0.5814 & 0.6033 \\
\end{longtable}\endgroup


\chapter{Fine-tuning replications and operating points}
The following tables reproduce the primary three operating points and all seed-specific routed results. They preserve the development routing rather than replacing it with the best test candidate. ``M'' denotes MobileCLIP2-S2 and ``S'' denotes SigLIP2-B/16-224 in compact tables. AP, F1, recall and FPR use the exact cohort/split described in Chapter~\ref{ch:training}; public sleep and historical unattended labels retain their provisional qualifications.
\begingroup\footnotesize
\setlength{\LTpre}{5pt}\setlength{\LTpost}{8pt}
\begin{longtable}{@{}>{\raggedright\arraybackslash}p{0.1200\linewidth}>{\raggedright\arraybackslash}p{0.1527\linewidth}>{\raggedright\arraybackslash}p{0.1418\linewidth}>{\raggedright\arraybackslash}p{0.0982\linewidth}>{\raggedright\arraybackslash}p{0.0764\linewidth}>{\raggedright\arraybackslash}p{0.0873\linewidth}>{\raggedright\arraybackslash}p{0.0764\linewidth}>{\raggedright\arraybackslash}p{0.0873\linewidth}@{}}
\caption[Finetuning all operating]{All primary routed operating points; historical/provisional cohort qualifications apply.}\label{tab:finetuning_all_operating}\\
\toprule
\textbf{Model} & \textbf{Target} & \textbf{Cutoff rule} & \textbf{Cutoff} & \textbf{PPV} & \textbf{Recall} & \textbf{F1} & \textbf{FPR} \\
\midrule\endfirsthead
\multicolumn{8}{l}{\textit{Table \thetable\ continued}}\\
\toprule
\textbf{Model} & \textbf{Target} & \textbf{Cutoff rule} & \textbf{Cutoff} & \textbf{PPV} & \textbf{Recall} & \textbf{F1} & \textbf{FPR} \\
\midrule\endhead
\midrule\multicolumn{8}{r}{\textit{Continued on next page}}\\\endfoot
\bottomrule\endlastfoot
Mobile\allowbreak{}CLIP & Smoking & raw & 0.000 & 0.0\% & 0.0\% & 0.000 & 1.2\% \\
Mobile\allowbreak{}CLIP & Smoking & dev max f1 & -4.416 & 60.7\% & 93.1\% & 0.735 & 93.7\% \\
Mobile\allowbreak{}CLIP & Smoking & dev fpr budget & -0.760 & 50.0\% & 3.1\% & 0.058 & 4.7\% \\
Mobile\allowbreak{}CLIP & Dumping / littering & raw & 0.000 & 78.8\% & 82.0\% & 0.804 & 22.0\% \\
Mobile\allowbreak{}CLIP & Dumping / littering & dev max f1 & -1.372 & 78.9\% & 90.0\% & 0.841 & 24.0\% \\
Mobile\allowbreak{}CLIP & Dumping / littering & dev fpr budget & 4.225 & 90.5\% & 38.0\% & 0.535 & 4.0\% \\
Mobile\allowbreak{}CLIP & Loitering & raw & 0.000 & 0.0\% & 0.0\% & 0.000 & 9.4\% \\
Mobile\allowbreak{}CLIP & Loitering & dev max f1 & -0.472 & 0.0\% & 0.0\% & 0.000 & 12.5\% \\
Mobile\allowbreak{}CLIP & Loitering & dev fpr budget & -0.472 & 0.0\% & 0.0\% & 0.000 & 12.5\% \\
Mobile\allowbreak{}CLIP & Public sleep & raw & 0.000 & 83.3\% & 83.3\% & 0.833 & 8.3\% \\
Mobile\allowbreak{}CLIP & Public sleep & dev max f1 & 0.779 & 100.0\% & 83.3\% & 0.909 & 0.0\% \\
Mobile\allowbreak{}CLIP & Public sleep & dev fpr budget & 0.779 & 100.0\% & 83.3\% & 0.909 & 0.0\% \\
Mobile\allowbreak{}CLIP & Unattended & raw & 0.000 & 100.0\% & 25.6\% & 0.408 & 0.0\% \\
Mobile\allowbreak{}CLIP & Unattended & dev max f1 & -5.249 & 25.6\% & 100.0\% & 0.408 & 97.6\% \\
Mobile\allowbreak{}CLIP & Unattended & dev fpr budget & -3.311 & 47.0\% & 96.3\% & 0.631 & 36.5\% \\
SigLIP & Smoking & raw & 0.000 & 68.3\% & 10.9\% & 0.189 & 7.9\% \\
SigLIP & Smoking & dev max f1 & -1.445 & 60.3\% & 91.3\% & 0.727 & 93.3\% \\
SigLIP & Smoking & dev fpr budget & 0.778 & 100.0\% & 0.3\% & 0.005 & 0.0\% \\
SigLIP & Dumping / littering & raw & 0.000 & 75.5\% & 80.0\% & 0.777 & 26.0\% \\
SigLIP & Dumping / littering & dev max f1 & -0.245 & 76.8\% & 86.0\% & 0.811 & 26.0\% \\
SigLIP & Dumping / littering & dev fpr budget & 1.790 & 85.7\% & 60.0\% & 0.706 & 10.0\% \\
SigLIP & Loitering & raw & 0.000 & 0.0\% & 0.0\% & 0.000 & 0.0\% \\
SigLIP & Loitering & dev max f1 & -2.409 & 2.2\% & 9.1\% & 0.035 & 70.3\% \\
SigLIP & Loitering & dev fpr budget & -1.556 & 0.0\% & 0.0\% & 0.000 & 12.5\% \\
SigLIP & Public sleep & raw & 0.000 & 85.7\% & 100.0\% & 0.923 & 8.3\% \\
SigLIP & Public sleep & dev max f1 & 2.163 & 100.0\% & 33.3\% & 0.500 & 0.0\% \\
SigLIP & Public sleep & dev fpr budget & 2.163 & 100.0\% & 33.3\% & 0.500 & 0.0\% \\
SigLIP & Unattended & raw & 0.000 & 96.7\% & 61.2\% & 0.749 & 0.7\% \\
SigLIP & Unattended & dev max f1 & -5.255 & 25.2\% & 100.0\% & 0.403 & 99.6\% \\
SigLIP & Unattended & dev fpr budget & -4.791 & 25.3\% & 100.0\% & 0.403 & 99.4\% \\
\end{longtable}\endgroup

\begingroup\footnotesize
\setlength{\LTpre}{5pt}\setlength{\LTpost}{8pt}
\begin{longtable}{@{}>{\raggedright\arraybackslash}p{0.1266\linewidth}>{\raggedright\arraybackslash}p{0.0575\linewidth}>{\raggedright\arraybackslash}p{0.1726\linewidth}>{\raggedright\arraybackslash}p{0.1496\linewidth}>{\raggedright\arraybackslash}p{0.0748\linewidth}>{\raggedright\arraybackslash}p{0.0748\linewidth}>{\raggedright\arraybackslash}p{0.0921\linewidth}>{\raggedright\arraybackslash}p{0.0921\linewidth}@{}}
\caption[Finetuning seeds]{All seed-specific routed results at maximum-dev-F1 cutoffs; seed 42 remains primary.}\label{tab:finetuning_seeds}\\
\toprule
\textbf{Model} & \textbf{Seed} & \textbf{Target} & \textbf{Dev route} & \textbf{AP} & \textbf{F1} & \textbf{Recall} & \textbf{FPR} \\
\midrule\endfirsthead
\multicolumn{8}{l}{\textit{Table \thetable\ continued}}\\
\toprule
\textbf{Model} & \textbf{Seed} & \textbf{Target} & \textbf{Dev route} & \textbf{AP} & \textbf{F1} & \textbf{Recall} & \textbf{FPR} \\
\midrule\endhead
\midrule\multicolumn{8}{r}{\textit{Continued on next page}}\\\endfoot
\bottomrule\endlastfoot
Mobile\allowbreak{}CLIP & 42 & Smoking & finetuned & 0.564 & 0.735 & 93.1\% & 93.7\% \\
Mobile\allowbreak{}CLIP & 42 & Dumping / littering & finetuned & 0.858 & 0.841 & 90.0\% & 24.0\% \\
Mobile\allowbreak{}CLIP & 42 & Loitering & finetuned & 0.093 & 0.000 & 0.0\% & 12.5\% \\
Mobile\allowbreak{}CLIP & 42 & Public sleep & finetuned & 0.976 & 0.909 & 83.3\% & 0.0\% \\
Mobile\allowbreak{}CLIP & 42 & Unattended & finetuned & 0.652 & 0.408 & 100.0\% & 97.6\% \\
Mobile\allowbreak{}CLIP & 43 & Smoking & linear probe & 0.594 & 0.752 & 98.0\% & 97.2\% \\
Mobile\allowbreak{}CLIP & 43 & Dumping / littering & linear probe & 0.861 & 0.847 & 94.0\% & 28.0\% \\
Mobile\allowbreak{}CLIP & 43 & Loitering & finetuned & 0.099 & 0.097 & 27.3\% & 75.0\% \\
Mobile\allowbreak{}CLIP & 43 & Public sleep & finetuned & 1.000 & 0.909 & 83.3\% & 0.0\% \\
Mobile\allowbreak{}CLIP & 43 & Unattended & finetuned & 0.697 & 0.405 & 100.0\% & 98.6\% \\
Mobile\allowbreak{}CLIP & 44 & Smoking & finetuned & 0.590 & 0.756 & 100.0\% & 100.0\% \\
Mobile\allowbreak{}CLIP & 44 & Dumping / littering & finetuned & 0.799 & 0.760 & 76.0\% & 24.0\% \\
Mobile\allowbreak{}CLIP & 44 & Loitering & finetuned & 0.111 & 0.000 & 0.0\% & 9.4\% \\
Mobile\allowbreak{}CLIP & 44 & Public sleep & finetuned & 0.976 & 0.909 & 83.3\% & 0.0\% \\
Mobile\allowbreak{}CLIP & 44 & Unattended & finetuned & 0.513 & 0.406 & 100.0\% & 98.5\% \\
SigLIP & 42 & Smoking & finetuned & 0.650 & 0.727 & 91.3\% & 93.3\% \\
SigLIP & 42 & Dumping / littering & linear probe & 0.885 & 0.811 & 86.0\% & 26.0\% \\
SigLIP & 42 & Loitering & zero shot & 0.097 & 0.035 & 9.1\% & 70.3\% \\
SigLIP & 42 & Public sleep & linear probe & 0.915 & 0.500 & 33.3\% & 0.0\% \\
SigLIP & 42 & Unattended & finetuned & 0.896 & 0.403 & 100.0\% & 99.6\% \\
SigLIP & 43 & Smoking & finetuned & 0.665 & 0.756 & 100.0\% & 100.0\% \\
SigLIP & 43 & Dumping / littering & linear probe & 0.885 & 0.811 & 86.0\% & 26.0\% \\
SigLIP & 43 & Loitering & zero shot & 0.097 & 0.035 & 9.1\% & 70.3\% \\
SigLIP & 43 & Public sleep & linear probe & 0.915 & 0.500 & 33.3\% & 0.0\% \\
SigLIP & 43 & Unattended & finetuned & 0.902 & 0.403 & 100.0\% & 99.4\% \\
SigLIP & 44 & Smoking & finetuned & 0.692 & 0.755 & 99.5\% & 99.6\% \\
SigLIP & 44 & Dumping / littering & linear probe & 0.885 & 0.811 & 86.0\% & 26.0\% \\
SigLIP & 44 & Loitering & finetuned & 0.095 & 0.000 & 0.0\% & 7.8\% \\
SigLIP & 44 & Public sleep & linear probe & 0.915 & 0.500 & 33.3\% & 0.0\% \\
SigLIP & 44 & Unattended & linear probe & 0.855 & 0.402 & 100.0\% & 100.0\% \\
\end{longtable}\endgroup


\chapter{All 330 strategy-development configurations}
This is a complete development table: 33 candidates for each of five targets and two backbones. It is not a test-selected leaderboard. Strategy names preserve interpolation coefficients and representation recipes; the ``setting'' column gives linear $C$ or an MLP selected epoch. The primary per-task choice maximizes mean development-source AP. Complete held-out family operating-point data remain in the original CSV and are included in the source package's numeric exports.
\begingroup\footnotesize
\setlength{\LTpre}{5pt}\setlength{\LTpost}{8pt}
\begin{longtable}{@{}>{\raggedright\arraybackslash}p{0.0575\linewidth}>{\raggedright\arraybackslash}p{0.1726\linewidth}>{\raggedright\arraybackslash}p{0.2532\linewidth}>{\raggedright\arraybackslash}p{0.0921\linewidth}>{\raggedright\arraybackslash}p{0.0690\linewidth}>{\raggedright\arraybackslash}p{0.0978\linewidth}>{\raggedright\arraybackslash}p{0.0978\linewidth}@{}}
\caption[Strategy all candidates]{All 330 strategy candidates: development scores only; M=MobileCLIP, S=SigLIP.}\label{tab:strategy_all_candidates}\\
\toprule
\textbf{M/S} & \textbf{Target} & \textbf{Representation} & \textbf{Head} & \textbf{C/ep} & \textbf{Mean AP} & \textbf{Worst AP} \\
\midrule\endfirsthead
\multicolumn{7}{l}{\textit{Table \thetable\ continued}}\\
\toprule
\textbf{M/S} & \textbf{Target} & \textbf{Representation} & \textbf{Head} & \textbf{C/ep} & \textbf{Mean AP} & \textbf{Worst AP} \\
\midrule\endhead
\midrule\multicolumn{7}{r}{\textit{Continued on next page}}\\\endfoot
\bottomrule\endlastfoot
M & Smoking & endpoint a0 & zero\_shot & -- & 0.4450 & 0.1929 \\
M & Smoking & full pool a0 & linear & 0.1 & 0.7467 & 0.3703 \\
M & Smoking & full pool a0 & linear & 1 & 0.7526 & 0.3786 \\
M & Smoking & full pool a0 & linear & 10 & 0.7520 & 0.3922 \\
M & Smoking & full pool a0.25 & linear & 0.1 & 0.7663 & 0.3981 \\
M & Smoking & full pool a0.25 & linear & 1 & 0.7679 & 0.4047 \\
M & Smoking & full pool a0.25 & linear & 10 & 0.7582 & 0.3964 \\
M & Smoking & full pool a0.5 & linear & 0.1 & 0.7947 & 0.4642 \\
M & Smoking & full pool a0.5 & linear & 1 & 0.7940 & 0.4705 \\
M & Smoking & full pool a0.5 & linear & 10 & 0.7744 & 0.4477 \\
M & Smoking & full pool a1 & linear & 0.1 & 0.7604 & 0.4125 \\
M & Smoking & full pool a1 & linear & 1 & 0.7597 & 0.4312 \\
M & Smoking & full pool a1 & linear & 10 & 0.7612 & 0.4638 \\
M & Smoking & endpoint a0 & linear & 0.1 & 0.7467 & 0.3703 \\
M & Smoking & endpoint a0 & linear & 1 & 0.7527 & 0.3788 \\
M & Smoking & endpoint a0 & linear & 10 & 0.7520 & 0.3922 \\
M & Smoking & full ordered a0 & linear & 0.1 & 0.7467 & 0.3703 \\
M & Smoking & full ordered a0 & linear & 1 & 0.7526 & 0.3786 \\
M & Smoking & full ordered a0 & linear & 10 & 0.7490 & 0.3909 \\
M & Smoking & full ordered a1 & linear & 0.1 & 0.7604 & 0.4125 \\
M & Smoking & full ordered a1 & linear & 1 & 0.7597 & 0.4312 \\
M & Smoking & full ordered a1 & linear & 10 & 0.7612 & 0.4638 \\
M & Smoking & full crop pool a0 & linear & 0.1 & 0.7650 & 0.4320 \\
M & Smoking & full crop pool a0 & linear & 1 & 0.7837 & 0.4725 \\
M & Smoking & full crop pool a0 & linear & 10 & 0.7999 & 0.5095 \\
M & Smoking & full crop ordered a0 & linear & 0.1 & 0.7650 & 0.4320 \\
M & Smoking & full crop ordered a0 & linear & 1 & 0.7837 & 0.4725 \\
M & Smoking & full crop ordered a0 & linear & 10 & 0.7977 & 0.5052 \\
M & Smoking & full geometry a0 & linear & 0.1 & 0.7501 & 0.3790 \\
M & Smoking & full geometry a0 & linear & 1 & 0.7499 & 0.3765 \\
M & Smoking & full geometry a0 & linear & 10 & 0.7499 & 0.3874 \\
M & Smoking & full pool a0 & mlp & 19 & 0.7455 & 0.3726 \\
M & Smoking & full crop pool a0 & mlp & 28 & 0.7933 & 0.4986 \\
M & Dumping / littering & endpoint a0 & zero\_shot & -- & 0.4671 & 0.4671 \\
M & Dumping / littering & full pool a0 & linear & 0.1 & 0.8491 & 0.8491 \\
M & Dumping / littering & full pool a0 & linear & 1 & 0.8462 & 0.8462 \\
M & Dumping / littering & full pool a0 & linear & 10 & 0.8880 & 0.8880 \\
M & Dumping / littering & full pool a0.25 & linear & 0.1 & 0.8864 & 0.8864 \\
M & Dumping / littering & full pool a0.25 & linear & 1 & 0.8986 & 0.8986 \\
M & Dumping / littering & full pool a0.25 & linear & 10 & 0.9103 & 0.9103 \\
M & Dumping / littering & full pool a0.5 & linear & 0.1 & 0.9231 & 0.9231 \\
M & Dumping / littering & full pool a0.5 & linear & 1 & 0.9185 & 0.9185 \\
M & Dumping / littering & full pool a0.5 & linear & 10 & 0.9184 & 0.9184 \\
M & Dumping / littering & full pool a1 & linear & 0.1 & 0.9123 & 0.9123 \\
M & Dumping / littering & full pool a1 & linear & 1 & 0.9096 & 0.9096 \\
M & Dumping / littering & full pool a1 & linear & 10 & 0.9112 & 0.9112 \\
M & Dumping / littering & endpoint a0 & linear & 0.1 & 0.8486 & 0.8486 \\
M & Dumping / littering & endpoint a0 & linear & 1 & 0.8659 & 0.8659 \\
M & Dumping / littering & endpoint a0 & linear & 10 & 0.8974 & 0.8974 \\
M & Dumping / littering & full ordered a0 & linear & 0.1 & 0.8337 & 0.8337 \\
M & Dumping / littering & full ordered a0 & linear & 1 & 0.8380 & 0.8380 \\
M & Dumping / littering & full ordered a0 & linear & 10 & 0.8839 & 0.8839 \\
M & Dumping / littering & full ordered a1 & linear & 0.1 & 0.9065 & 0.9065 \\
M & Dumping / littering & full ordered a1 & linear & 1 & 0.9134 & 0.9134 \\
M & Dumping / littering & full ordered a1 & linear & 10 & 0.9175 & 0.9175 \\
M & Dumping / littering & full crop pool a0 & linear & 0.1 & 0.8241 & 0.8241 \\
M & Dumping / littering & full crop pool a0 & linear & 1 & 0.8609 & 0.8609 \\
M & Dumping / littering & full crop pool a0 & linear & 10 & 0.9062 & 0.9062 \\
M & Dumping / littering & full crop ordered a0 & linear & 0.1 & 0.8201 & 0.8201 \\
M & Dumping / littering & full crop ordered a0 & linear & 1 & 0.8400 & 0.8400 \\
M & Dumping / littering & full crop ordered a0 & linear & 10 & 0.8668 & 0.8668 \\
M & Dumping / littering & full geometry a0 & linear & 0.1 & 0.7218 & 0.7218 \\
M & Dumping / littering & full geometry a0 & linear & 1 & 0.8470 & 0.8470 \\
M & Dumping / littering & full geometry a0 & linear & 10 & 0.9153 & 0.9153 \\
M & Dumping / littering & full pool a0 & mlp & 1 & 0.8556 & 0.8556 \\
M & Dumping / littering & full crop pool a0 & mlp & 12 & 0.8415 & 0.8415 \\
M & Loitering & endpoint a0 & zero\_shot & -- & 0.1747 & 0.1747 \\
M & Loitering & full pool a0 & linear & 0.1 & 0.1058 & 0.1058 \\
M & Loitering & full pool a0 & linear & 1 & 0.1096 & 0.1096 \\
M & Loitering & full pool a0 & linear & 10 & 0.1199 & 0.1199 \\
M & Loitering & full pool a0.25 & linear & 0.1 & 0.1340 & 0.1340 \\
M & Loitering & full pool a0.25 & linear & 1 & 0.1429 & 0.1429 \\
M & Loitering & full pool a0.25 & linear & 10 & 0.1720 & 0.1720 \\
M & Loitering & full pool a0.5 & linear & 0.1 & 0.1926 & 0.1926 \\
M & Loitering & full pool a0.5 & linear & 1 & 0.2216 & 0.2216 \\
M & Loitering & full pool a0.5 & linear & 10 & 0.2251 & 0.2251 \\
M & Loitering & full pool a1 & linear & 0.1 & 0.2915 & 0.2915 \\
M & Loitering & full pool a1 & linear & 1 & 0.2792 & 0.2792 \\
M & Loitering & full pool a1 & linear & 10 & 0.2620 & 0.2620 \\
M & Loitering & endpoint a0 & linear & 0.1 & 0.1183 & 0.1183 \\
M & Loitering & endpoint a0 & linear & 1 & 0.1254 & 0.1254 \\
M & Loitering & endpoint a0 & linear & 10 & 0.1374 & 0.1374 \\
M & Loitering & full ordered a0 & linear & 0.1 & 0.1156 & 0.1156 \\
M & Loitering & full ordered a0 & linear & 1 & 0.1216 & 0.1216 \\
M & Loitering & full ordered a0 & linear & 10 & 0.1377 & 0.1377 \\
M & Loitering & full ordered a1 & linear & 0.1 & 0.2332 & 0.2332 \\
M & Loitering & full ordered a1 & linear & 1 & 0.2346 & 0.2346 \\
M & Loitering & full ordered a1 & linear & 10 & 0.2105 & 0.2105 \\
M & Loitering & full crop pool a0 & linear & 0.1 & 0.2170 & 0.2170 \\
M & Loitering & full crop pool a0 & linear & 1 & 0.1972 & 0.1972 \\
M & Loitering & full crop pool a0 & linear & 10 & 0.1884 & 0.1884 \\
M & Loitering & full crop ordered a0 & linear & 0.1 & 0.2410 & 0.2410 \\
M & Loitering & full crop ordered a0 & linear & 1 & 0.2018 & 0.2018 \\
M & Loitering & full crop ordered a0 & linear & 10 & 0.1777 & 0.1777 \\
M & Loitering & full geometry a0 & linear & 0.1 & 0.4858 & 0.4858 \\
M & Loitering & full geometry a0 & linear & 1 & 0.3374 & 0.3374 \\
M & Loitering & full geometry a0 & linear & 10 & 0.2079 & 0.2079 \\
M & Loitering & full pool a0 & mlp & 5 & 0.1052 & 0.1052 \\
M & Loitering & full crop pool a0 & mlp & 5 & 0.1913 & 0.1913 \\
M & Public sleep & endpoint a0 & zero\_shot & -- & 0.5420 & 0.5420 \\
M & Public sleep & full pool a0 & linear & 0.1 & 0.9683 & 0.9683 \\
M & Public sleep & full pool a0 & linear & 1 & 0.9683 & 0.9683 \\
M & Public sleep & full pool a0 & linear & 10 & 0.9683 & 0.9683 \\
M & Public sleep & full pool a0.25 & linear & 0.1 & 1.0000 & 1.0000 \\
M & Public sleep & full pool a0.25 & linear & 1 & 1.0000 & 1.0000 \\
M & Public sleep & full pool a0.25 & linear & 10 & 1.0000 & 1.0000 \\
M & Public sleep & full pool a0.5 & linear & 0.1 & 1.0000 & 1.0000 \\
M & Public sleep & full pool a0.5 & linear & 1 & 1.0000 & 1.0000 \\
M & Public sleep & full pool a0.5 & linear & 10 & 1.0000 & 1.0000 \\
M & Public sleep & full pool a1 & linear & 0.1 & 1.0000 & 1.0000 \\
M & Public sleep & full pool a1 & linear & 1 & 1.0000 & 1.0000 \\
M & Public sleep & full pool a1 & linear & 10 & 1.0000 & 1.0000 \\
M & Public sleep & endpoint a0 & linear & 0.1 & 0.9683 & 0.9683 \\
M & Public sleep & endpoint a0 & linear & 1 & 0.9683 & 0.9683 \\
M & Public sleep & endpoint a0 & linear & 10 & 0.9683 & 0.9683 \\
M & Public sleep & full ordered a0 & linear & 0.1 & 0.9683 & 0.9683 \\
M & Public sleep & full ordered a0 & linear & 1 & 0.9683 & 0.9683 \\
M & Public sleep & full ordered a0 & linear & 10 & 0.9683 & 0.9683 \\
M & Public sleep & full ordered a1 & linear & 0.1 & 1.0000 & 1.0000 \\
M & Public sleep & full ordered a1 & linear & 1 & 1.0000 & 1.0000 \\
M & Public sleep & full ordered a1 & linear & 10 & 1.0000 & 1.0000 \\
M & Public sleep & full crop pool a0 & linear & 0.1 & 0.9683 & 0.9683 \\
M & Public sleep & full crop pool a0 & linear & 1 & 0.9821 & 0.9821 \\
M & Public sleep & full crop pool a0 & linear & 10 & 0.9821 & 0.9821 \\
M & Public sleep & full crop ordered a0 & linear & 0.1 & 0.9683 & 0.9683 \\
M & Public sleep & full crop ordered a0 & linear & 1 & 0.9821 & 0.9821 \\
M & Public sleep & full crop ordered a0 & linear & 10 & 0.9821 & 0.9821 \\
M & Public sleep & full geometry a0 & linear & 0.1 & 0.8955 & 0.8955 \\
M & Public sleep & full geometry a0 & linear & 1 & 0.9240 & 0.9240 \\
M & Public sleep & full geometry a0 & linear & 10 & 0.9478 & 0.9478 \\
M & Public sleep & full pool a0 & mlp & 1 & 1.0000 & 1.0000 \\
M & Public sleep & full crop pool a0 & mlp & 1 & 1.0000 & 1.0000 \\
M & Unattended & endpoint a0 & zero\_shot & -- & 0.4967 & 0.4967 \\
M & Unattended & full pool a0 & linear & 0.1 & 0.6441 & 0.6441 \\
M & Unattended & full pool a0 & linear & 1 & 0.6221 & 0.6221 \\
M & Unattended & full pool a0 & linear & 10 & 0.5933 & 0.5933 \\
M & Unattended & full pool a0.25 & linear & 0.1 & 0.7052 & 0.7052 \\
M & Unattended & full pool a0.25 & linear & 1 & 0.7207 & 0.7207 \\
M & Unattended & full pool a0.25 & linear & 10 & 0.6968 & 0.6968 \\
M & Unattended & full pool a0.5 & linear & 0.1 & 0.7797 & 0.7797 \\
M & Unattended & full pool a0.5 & linear & 1 & 0.7714 & 0.7714 \\
M & Unattended & full pool a0.5 & linear & 10 & 0.7633 & 0.7633 \\
M & Unattended & full pool a1 & linear & 0.1 & 0.8310 & 0.8310 \\
M & Unattended & full pool a1 & linear & 1 & 0.8477 & 0.8477 \\
M & Unattended & full pool a1 & linear & 10 & 0.8641 & 0.8641 \\
M & Unattended & endpoint a0 & linear & 0.1 & 0.6711 & 0.6711 \\
M & Unattended & endpoint a0 & linear & 1 & 0.6275 & 0.6275 \\
M & Unattended & endpoint a0 & linear & 10 & 0.7049 & 0.7049 \\
M & Unattended & full ordered a0 & linear & 0.1 & 0.5108 & 0.5108 \\
M & Unattended & full ordered a0 & linear & 1 & 0.5405 & 0.5405 \\
M & Unattended & full ordered a0 & linear & 10 & 0.6042 & 0.6042 \\
M & Unattended & full ordered a1 & linear & 0.1 & 0.9222 & 0.9222 \\
M & Unattended & full ordered a1 & linear & 1 & 0.9442 & 0.9442 \\
M & Unattended & full ordered a1 & linear & 10 & 0.9469 & 0.9469 \\
M & Unattended & full crop pool a0 & linear & 0.1 & 0.7582 & 0.7582 \\
M & Unattended & full crop pool a0 & linear & 1 & 0.7826 & 0.7826 \\
M & Unattended & full crop pool a0 & linear & 10 & 0.5570 & 0.5570 \\
M & Unattended & full crop ordered a0 & linear & 0.1 & 0.7149 & 0.7149 \\
M & Unattended & full crop ordered a0 & linear & 1 & 0.8159 & 0.8159 \\
M & Unattended & full crop ordered a0 & linear & 10 & 0.5884 & 0.5884 \\
M & Unattended & full geometry a0 & linear & 0.1 & 0.6529 & 0.6529 \\
M & Unattended & full geometry a0 & linear & 1 & 0.6987 & 0.6987 \\
M & Unattended & full geometry a0 & linear & 10 & 0.6460 & 0.6460 \\
M & Unattended & full pool a0 & mlp & 5 & 0.6525 & 0.6525 \\
M & Unattended & full crop pool a0 & mlp & 14 & 0.7569 & 0.7569 \\
S & Smoking & endpoint a0 & zero\_shot & -- & 0.5667 & 0.3957 \\
S & Smoking & full pool a0 & linear & 0.1 & 0.7509 & 0.3737 \\
S & Smoking & full pool a0 & linear & 1 & 0.7668 & 0.3705 \\
S & Smoking & full pool a0 & linear & 10 & 0.7770 & 0.3857 \\
S & Smoking & full pool a0.25 & linear & 0.1 & 0.7712 & 0.3841 \\
S & Smoking & full pool a0.25 & linear & 1 & 0.7775 & 0.3793 \\
S & Smoking & full pool a0.25 & linear & 10 & 0.7850 & 0.3943 \\
S & Smoking & full pool a0.5 & linear & 0.1 & 0.7825 & 0.3955 \\
S & Smoking & full pool a0.5 & linear & 1 & 0.7838 & 0.3891 \\
S & Smoking & full pool a0.5 & linear & 10 & 0.7813 & 0.3848 \\
S & Smoking & full pool a1 & linear & 0.1 & 0.7951 & 0.4105 \\
S & Smoking & full pool a1 & linear & 1 & 0.7920 & 0.3984 \\
S & Smoking & full pool a1 & linear & 10 & 0.7828 & 0.3828 \\
S & Smoking & endpoint a0 & linear & 0.1 & 0.7509 & 0.3737 \\
S & Smoking & endpoint a0 & linear & 1 & 0.7668 & 0.3705 \\
S & Smoking & endpoint a0 & linear & 10 & 0.7783 & 0.3919 \\
S & Smoking & full ordered a0 & linear & 0.1 & 0.7510 & 0.3727 \\
S & Smoking & full ordered a0 & linear & 1 & 0.7668 & 0.3705 \\
S & Smoking & full ordered a0 & linear & 10 & 0.7816 & 0.3902 \\
S & Smoking & full ordered a1 & linear & 0.1 & 0.7951 & 0.4105 \\
S & Smoking & full ordered a1 & linear & 1 & 0.7920 & 0.3984 \\
S & Smoking & full ordered a1 & linear & 10 & 0.7828 & 0.3828 \\
S & Smoking & full crop pool a0 & linear & 0.1 & 0.7721 & 0.4271 \\
S & Smoking & full crop pool a0 & linear & 1 & 0.7966 & 0.4517 \\
S & Smoking & full crop pool a0 & linear & 10 & 0.8124 & 0.4863 \\
S & Smoking & full crop ordered a0 & linear & 0.1 & 0.7721 & 0.4271 \\
S & Smoking & full crop ordered a0 & linear & 1 & 0.7966 & 0.4517 \\
S & Smoking & full crop ordered a0 & linear & 10 & 0.8028 & 0.4629 \\
S & Smoking & full geometry a0 & linear & 0.1 & 0.7564 & 0.3906 \\
S & Smoking & full geometry a0 & linear & 1 & 0.7652 & 0.3714 \\
S & Smoking & full geometry a0 & linear & 10 & 0.7757 & 0.3789 \\
S & Smoking & full pool a0 & mlp & 31 & 0.7777 & 0.3714 \\
S & Smoking & full crop pool a0 & mlp & 39 & 0.8209 & 0.5053 \\
S & Dumping / littering & endpoint a0 & zero\_shot & -- & 0.6288 & 0.6288 \\
S & Dumping / littering & full pool a0 & linear & 0.1 & 0.8704 & 0.8704 \\
S & Dumping / littering & full pool a0 & linear & 1 & 0.8846 & 0.8846 \\
S & Dumping / littering & full pool a0 & linear & 10 & 0.9185 & 0.9185 \\
S & Dumping / littering & full pool a0.25 & linear & 0.1 & 0.8789 & 0.8789 \\
S & Dumping / littering & full pool a0.25 & linear & 1 & 0.8858 & 0.8858 \\
S & Dumping / littering & full pool a0.25 & linear & 10 & 0.9207 & 0.9207 \\
S & Dumping / littering & full pool a0.5 & linear & 0.1 & 0.8861 & 0.8861 \\
S & Dumping / littering & full pool a0.5 & linear & 1 & 0.8954 & 0.8954 \\
S & Dumping / littering & full pool a0.5 & linear & 10 & 0.9146 & 0.9146 \\
S & Dumping / littering & full pool a1 & linear & 0.1 & 0.8823 & 0.8823 \\
S & Dumping / littering & full pool a1 & linear & 1 & 0.8966 & 0.8966 \\
S & Dumping / littering & full pool a1 & linear & 10 & 0.9010 & 0.9010 \\
S & Dumping / littering & endpoint a0 & linear & 0.1 & 0.8608 & 0.8608 \\
S & Dumping / littering & endpoint a0 & linear & 1 & 0.8763 & 0.8763 \\
S & Dumping / littering & endpoint a0 & linear & 10 & 0.9145 & 0.9145 \\
S & Dumping / littering & full ordered a0 & linear & 0.1 & 0.8552 & 0.8552 \\
S & Dumping / littering & full ordered a0 & linear & 1 & 0.8692 & 0.8692 \\
S & Dumping / littering & full ordered a0 & linear & 10 & 0.9059 & 0.9059 \\
S & Dumping / littering & full ordered a1 & linear & 0.1 & 0.8848 & 0.8848 \\
S & Dumping / littering & full ordered a1 & linear & 1 & 0.8880 & 0.8880 \\
S & Dumping / littering & full ordered a1 & linear & 10 & 0.9060 & 0.9060 \\
S & Dumping / littering & full crop pool a0 & linear & 0.1 & 0.8275 & 0.8275 \\
S & Dumping / littering & full crop pool a0 & linear & 1 & 0.8581 & 0.8581 \\
S & Dumping / littering & full crop pool a0 & linear & 10 & 0.9017 & 0.9017 \\
S & Dumping / littering & full crop ordered a0 & linear & 0.1 & 0.8399 & 0.8399 \\
S & Dumping / littering & full crop ordered a0 & linear & 1 & 0.8608 & 0.8608 \\
S & Dumping / littering & full crop ordered a0 & linear & 10 & 0.8996 & 0.8996 \\
S & Dumping / littering & full geometry a0 & linear & 0.1 & 0.6962 & 0.6962 \\
S & Dumping / littering & full geometry a0 & linear & 1 & 0.8542 & 0.8542 \\
S & Dumping / littering & full geometry a0 & linear & 10 & 0.9263 & 0.9263 \\
S & Dumping / littering & full pool a0 & mlp & 13 & 0.8651 & 0.8651 \\
S & Dumping / littering & full crop pool a0 & mlp & 12 & 0.8122 & 0.8122 \\
S & Loitering & endpoint a0 & zero\_shot & -- & 0.2290 & 0.2290 \\
S & Loitering & full pool a0 & linear & 0.1 & 0.0845 & 0.0845 \\
S & Loitering & full pool a0 & linear & 1 & 0.0860 & 0.0860 \\
S & Loitering & full pool a0 & linear & 10 & 0.0948 & 0.0948 \\
S & Loitering & full pool a0.25 & linear & 0.1 & 0.0889 & 0.0889 \\
S & Loitering & full pool a0.25 & linear & 1 & 0.0935 & 0.0935 \\
S & Loitering & full pool a0.25 & linear & 10 & 0.1044 & 0.1044 \\
S & Loitering & full pool a0.5 & linear & 0.1 & 0.0968 & 0.0968 \\
S & Loitering & full pool a0.5 & linear & 1 & 0.1014 & 0.1014 \\
S & Loitering & full pool a0.5 & linear & 10 & 0.1228 & 0.1228 \\
S & Loitering & full pool a1 & linear & 0.1 & 0.1144 & 0.1144 \\
S & Loitering & full pool a1 & linear & 1 & 0.1293 & 0.1293 \\
S & Loitering & full pool a1 & linear & 10 & 0.1514 & 0.1514 \\
S & Loitering & endpoint a0 & linear & 0.1 & 0.0926 & 0.0926 \\
S & Loitering & endpoint a0 & linear & 1 & 0.0934 & 0.0934 \\
S & Loitering & endpoint a0 & linear & 10 & 0.0986 & 0.0986 \\
S & Loitering & full ordered a0 & linear & 0.1 & 0.0842 & 0.0842 \\
S & Loitering & full ordered a0 & linear & 1 & 0.0848 & 0.0848 \\
S & Loitering & full ordered a0 & linear & 10 & 0.0896 & 0.0896 \\
S & Loitering & full ordered a1 & linear & 0.1 & 0.1045 & 0.1045 \\
S & Loitering & full ordered a1 & linear & 1 & 0.1167 & 0.1167 \\
S & Loitering & full ordered a1 & linear & 10 & 0.1399 & 0.1399 \\
S & Loitering & full crop pool a0 & linear & 0.1 & 0.2598 & 0.2598 \\
S & Loitering & full crop pool a0 & linear & 1 & 0.2129 & 0.2129 \\
S & Loitering & full crop pool a0 & linear & 10 & 0.1942 & 0.1942 \\
S & Loitering & full crop ordered a0 & linear & 0.1 & 0.2785 & 0.2785 \\
S & Loitering & full crop ordered a0 & linear & 1 & 0.2206 & 0.2206 \\
S & Loitering & full crop ordered a0 & linear & 10 & 0.1919 & 0.1919 \\
S & Loitering & full geometry a0 & linear & 0.1 & 0.5223 & 0.5223 \\
S & Loitering & full geometry a0 & linear & 1 & 0.5063 & 0.5063 \\
S & Loitering & full geometry a0 & linear & 10 & 0.3094 & 0.3094 \\
S & Loitering & full pool a0 & mlp & 8 & 0.0871 & 0.0871 \\
S & Loitering & full crop pool a0 & mlp & 15 & 0.2120 & 0.2120 \\
S & Public sleep & endpoint a0 & zero\_shot & -- & 0.8524 & 0.8524 \\
S & Public sleep & full pool a0 & linear & 0.1 & 1.0000 & 1.0000 \\
S & Public sleep & full pool a0 & linear & 1 & 1.0000 & 1.0000 \\
S & Public sleep & full pool a0 & linear & 10 & 1.0000 & 1.0000 \\
S & Public sleep & full pool a0.25 & linear & 0.1 & 1.0000 & 1.0000 \\
S & Public sleep & full pool a0.25 & linear & 1 & 1.0000 & 1.0000 \\
S & Public sleep & full pool a0.25 & linear & 10 & 1.0000 & 1.0000 \\
S & Public sleep & full pool a0.5 & linear & 0.1 & 1.0000 & 1.0000 \\
S & Public sleep & full pool a0.5 & linear & 1 & 1.0000 & 1.0000 \\
S & Public sleep & full pool a0.5 & linear & 10 & 1.0000 & 1.0000 \\
S & Public sleep & full pool a1 & linear & 0.1 & 1.0000 & 1.0000 \\
S & Public sleep & full pool a1 & linear & 1 & 1.0000 & 1.0000 \\
S & Public sleep & full pool a1 & linear & 10 & 1.0000 & 1.0000 \\
S & Public sleep & endpoint a0 & linear & 0.1 & 1.0000 & 1.0000 \\
S & Public sleep & endpoint a0 & linear & 1 & 1.0000 & 1.0000 \\
S & Public sleep & endpoint a0 & linear & 10 & 1.0000 & 1.0000 \\
S & Public sleep & full ordered a0 & linear & 0.1 & 1.0000 & 1.0000 \\
S & Public sleep & full ordered a0 & linear & 1 & 1.0000 & 1.0000 \\
S & Public sleep & full ordered a0 & linear & 10 & 1.0000 & 1.0000 \\
S & Public sleep & full ordered a1 & linear & 0.1 & 1.0000 & 1.0000 \\
S & Public sleep & full ordered a1 & linear & 1 & 1.0000 & 1.0000 \\
S & Public sleep & full ordered a1 & linear & 10 & 1.0000 & 1.0000 \\
S & Public sleep & full crop pool a0 & linear & 0.1 & 1.0000 & 1.0000 \\
S & Public sleep & full crop pool a0 & linear & 1 & 1.0000 & 1.0000 \\
S & Public sleep & full crop pool a0 & linear & 10 & 1.0000 & 1.0000 \\
S & Public sleep & full crop ordered a0 & linear & 0.1 & 1.0000 & 1.0000 \\
S & Public sleep & full crop ordered a0 & linear & 1 & 1.0000 & 1.0000 \\
S & Public sleep & full crop ordered a0 & linear & 10 & 1.0000 & 1.0000 \\
S & Public sleep & full geometry a0 & linear & 0.1 & 0.9094 & 0.9094 \\
S & Public sleep & full geometry a0 & linear & 1 & 0.9617 & 0.9617 \\
S & Public sleep & full geometry a0 & linear & 10 & 0.9617 & 0.9617 \\
S & Public sleep & full pool a0 & mlp & 2 & 1.0000 & 1.0000 \\
S & Public sleep & full crop pool a0 & mlp & 1 & 1.0000 & 1.0000 \\
S & Unattended & endpoint a0 & zero\_shot & -- & 0.4672 & 0.4672 \\
S & Unattended & full pool a0 & linear & 0.1 & 0.6561 & 0.6561 \\
S & Unattended & full pool a0 & linear & 1 & 0.7002 & 0.7002 \\
S & Unattended & full pool a0 & linear & 10 & 0.7779 & 0.7779 \\
S & Unattended & full pool a0.25 & linear & 0.1 & 0.6751 & 0.6751 \\
S & Unattended & full pool a0.25 & linear & 1 & 0.7250 & 0.7250 \\
S & Unattended & full pool a0.25 & linear & 10 & 0.8025 & 0.8025 \\
S & Unattended & full pool a0.5 & linear & 0.1 & 0.6905 & 0.6905 \\
S & Unattended & full pool a0.5 & linear & 1 & 0.7498 & 0.7498 \\
S & Unattended & full pool a0.5 & linear & 10 & 0.8216 & 0.8216 \\
S & Unattended & full pool a1 & linear & 0.1 & 0.7272 & 0.7272 \\
S & Unattended & full pool a1 & linear & 1 & 0.8065 & 0.8065 \\
S & Unattended & full pool a1 & linear & 10 & 0.8376 & 0.8376 \\
S & Unattended & endpoint a0 & linear & 0.1 & 0.6085 & 0.6085 \\
S & Unattended & endpoint a0 & linear & 1 & 0.7093 & 0.7093 \\
S & Unattended & endpoint a0 & linear & 10 & 0.8502 & 0.8502 \\
S & Unattended & full ordered a0 & linear & 0.1 & 0.5221 & 0.5221 \\
S & Unattended & full ordered a0 & linear & 1 & 0.6074 & 0.6074 \\
S & Unattended & full ordered a0 & linear & 10 & 0.8534 & 0.8534 \\
S & Unattended & full ordered a1 & linear & 0.1 & 0.7060 & 0.7060 \\
S & Unattended & full ordered a1 & linear & 1 & 0.8496 & 0.8496 \\
S & Unattended & full ordered a1 & linear & 10 & 0.9681 & 0.9681 \\
S & Unattended & full crop pool a0 & linear & 0.1 & 0.5440 & 0.5440 \\
S & Unattended & full crop pool a0 & linear & 1 & 0.6116 & 0.6116 \\
S & Unattended & full crop pool a0 & linear & 10 & 0.6748 & 0.6748 \\
S & Unattended & full crop ordered a0 & linear & 0.1 & 0.4900 & 0.4900 \\
S & Unattended & full crop ordered a0 & linear & 1 & 0.5747 & 0.5747 \\
S & Unattended & full crop ordered a0 & linear & 10 & 0.7216 & 0.7216 \\
S & Unattended & full geometry a0 & linear & 0.1 & 0.6319 & 0.6319 \\
S & Unattended & full geometry a0 & linear & 1 & 0.7142 & 0.7142 \\
S & Unattended & full geometry a0 & linear & 10 & 0.7430 & 0.7430 \\
S & Unattended & full pool a0 & mlp & 3 & 0.6707 & 0.6707 \\
S & Unattended & full crop pool a0 & mlp & 51 & 0.6430 & 0.6430 \\
\end{longtable}\endgroup


\chapter{All 50 final video-pilot heads}
Final source-guarded heads and their epoch/seed, worst-source AP, mean-source AP and selected-route flag are reproduced below. Earlier prefix-split fits are diagnostic only. A hard-negative checkpoint selected before epoch four is not evidence of mining benefit. These are frozen-pretrained-encoder heads, distinct from the earlier six actual visual fine-tunes.
\begingroup\footnotesize
\setlength{\LTpre}{5pt}\setlength{\LTpost}{8pt}
\begin{longtable}{@{}>{\raggedright\arraybackslash}p{0.0506\linewidth}>{\raggedright\arraybackslash}p{0.1214\linewidth}>{\raggedright\arraybackslash}p{0.3542\linewidth}>{\raggedright\arraybackslash}p{0.0506\linewidth}>{\raggedright\arraybackslash}p{0.0607\linewidth}>{\raggedright\arraybackslash}p{0.0810\linewidth}>{\raggedright\arraybackslash}p{0.0810\linewidth}>{\raggedright\arraybackslash}p{0.0405\linewidth}@{}}
\caption[Pilot all heads]{All final source-guarded pilot heads; asterisk marks the primary development route.}\label{tab:pilot_all_heads}\\
\toprule
\textbf{M/S} & \textbf{Target} & \textbf{Head} & \textbf{Seed} & \textbf{Ep.} & \textbf{Worst AP} & \textbf{Mean AP} & \textbf{*} \\
\midrule\endfirsthead
\multicolumn{8}{l}{\textit{Table \thetable\ continued}}\\
\toprule
\textbf{M/S} & \textbf{Target} & \textbf{Head} & \textbf{Seed} & \textbf{Ep.} & \textbf{Worst AP} & \textbf{Mean AP} & \textbf{*} \\
\midrule\endhead
\midrule\multicolumn{8}{r}{\textit{Continued on next page}}\\\endfoot
\bottomrule\endlastfoot
M & Smoking & scene linear & -- & -- & 0.3733 & 0.6659 &  \\
M & Smoking & region linear & -- & -- & 0.4544 & 0.7368 &  \\
M & Smoking & region temporal seed 42 & 42 & 15 & 0.5427 & 0.7926 &  \\
M & Smoking & region temporal seed 43 & 43 & 22 & 0.5162 & 0.7723 &  \\
M & Smoking & region temporal seed 44 & 44 & 15 & 0.5173 & 0.7709 &  \\
M & Smoking & region track temporal seed 42 & 42 & 15 & 0.5492 & 0.7862 &  \\
M & Smoking & region track temporal seed 43 & 43 & 14 & 0.5274 & 0.7755 &  \\
M & Smoking & region track temporal seed 44 & 44 & 21 & 0.5175 & 0.7922 &  \\
M & Smoking & person supervised seed 42 & 42 & 60 & 0.5704 & 0.7967 &  \\
M & Smoking & person supervised seed 43 & 43 & 60 & 0.5723 & 0.7978 &  \\
M & Smoking & person supervised seed 44 & 44 & 60 & 0.5723 & 0.7994 &  \\
M & Smoking & person supervised hardneg seed 42 & 42 & 60 & 0.5736 & 0.7975 & * \\
M & Smoking & person supervised hardneg seed 43 & 43 & 60 & 0.5670 & 0.8004 &  \\
M & Smoking & person supervised hardneg seed 44 & 44 & 60 & 0.5719 & 0.8035 &  \\
M & Unattended & scene linear & -- & -- & 0.8130 & 0.8130 &  \\
M & Unattended & region linear & -- & -- & 0.8782 & 0.8782 & * \\
M & Unattended & region temporal seed 42 & 42 & 1 & 0.8146 & 0.8146 &  \\
M & Unattended & region temporal seed 43 & 43 & 4 & 0.7552 & 0.7552 &  \\
M & Unattended & region temporal seed 44 & 44 & 1 & 0.8316 & 0.8316 &  \\
M & Unattended & region track temporal seed 42 & 42 & 1 & 0.8259 & 0.8259 &  \\
M & Unattended & region track temporal seed 43 & 43 & 15 & 0.8297 & 0.8297 &  \\
M & Unattended & region track temporal seed 44 & 44 & 1 & 0.8318 & 0.8318 &  \\
M & Unattended & grounded track temporal seed 42 & 42 & 1 & 0.8067 & 0.8067 &  \\
M & Unattended & grounded track temporal seed 43 & 43 & 15 & 0.8377 & 0.8377 &  \\
M & Unattended & grounded track temporal seed 44 & 44 & 8 & 0.7828 & 0.7828 &  \\
S & Smoking & scene linear & -- & -- & 0.3723 & 0.6987 &  \\
S & Smoking & region linear & -- & -- & 0.4293 & 0.7264 &  \\
S & Smoking & region temporal seed 42 & 42 & 59 & 0.5810 & 0.7890 & * \\
S & Smoking & region temporal seed 43 & 43 & 23 & 0.5393 & 0.7818 &  \\
S & Smoking & region temporal seed 44 & 44 & 24 & 0.5435 & 0.7995 &  \\
S & Smoking & region track temporal seed 42 & 42 & 43 & 0.5710 & 0.7765 &  \\
S & Smoking & region track temporal seed 43 & 43 & 24 & 0.5530 & 0.7980 &  \\
S & Smoking & region track temporal seed 44 & 44 & 47 & 0.5774 & 0.7776 &  \\
S & Smoking & person supervised seed 42 & 42 & 9 & 0.4921 & 0.8164 &  \\
S & Smoking & person supervised seed 43 & 43 & 59 & 0.5531 & 0.8369 &  \\
S & Smoking & person supervised seed 44 & 44 & 60 & 0.5542 & 0.8372 &  \\
S & Smoking & person supervised hardneg seed 42 & 42 & 2 & 0.4890 & 0.8146 &  \\
S & Smoking & person supervised hardneg seed 43 & 43 & 60 & 0.5471 & 0.8348 &  \\
S & Smoking & person supervised hardneg seed 44 & 44 & 60 & 0.5506 & 0.8339 &  \\
S & Unattended & scene linear & -- & -- & 0.7520 & 0.7520 &  \\
S & Unattended & region linear & -- & -- & 0.7576 & 0.7576 &  \\
S & Unattended & region temporal seed 42 & 42 & 21 & 0.9700 & 0.9700 &  \\
S & Unattended & region temporal seed 43 & 43 & 39 & 0.9643 & 0.9643 &  \\
S & Unattended & region temporal seed 44 & 44 & 12 & 0.9477 & 0.9477 &  \\
S & Unattended & region track temporal seed 42 & 42 & 21 & 0.9552 & 0.9552 &  \\
S & Unattended & region track temporal seed 43 & 43 & 35 & 0.9710 & 0.9710 &  \\
S & Unattended & region track temporal seed 44 & 44 & 19 & 0.9120 & 0.9120 &  \\
S & Unattended & grounded track temporal seed 42 & 42 & 20 & 0.9730 & 0.9730 & * \\
S & Unattended & grounded track temporal seed 43 & 43 & 35 & 0.9701 & 0.9701 &  \\
S & Unattended & grounded track temporal seed 44 & 44 & 2 & 0.8805 & 0.8805 &  \\
\end{longtable}\endgroup


\chapter{All 38 latest smoking policies}
These policies reuse the existing saved heads and fixed equal sigmoid averages. Each cutoff satisfies the empirical development-source constraint; values are not calibrated probabilities. The selected flag is determined from worst/mean development recall with AP ties. A separate AP-selected alternative remains in the source JSON. All relevant held-out comparisons are retrospective.
\begingroup\footnotesize
\setlength{\LTpre}{5pt}\setlength{\LTpost}{8pt}
\begin{longtable}{@{}>{\raggedright\arraybackslash}p{0.0477\linewidth}>{\raggedright\arraybackslash}p{0.3341\linewidth}>{\raggedright\arraybackslash}p{0.0573\linewidth}>{\raggedright\arraybackslash}p{0.0859\linewidth}>{\raggedright\arraybackslash}p{0.0955\linewidth}>{\raggedright\arraybackslash}p{0.0955\linewidth}>{\raggedright\arraybackslash}p{0.0859\linewidth}>{\raggedright\arraybackslash}p{0.0382\linewidth}@{}}
\caption[Latest all policies]{All 38 smoking-policy development summaries; equal sigmoid averages are not calibrated probabilities.}\label{tab:latest_all_policies}\\
\toprule
\textbf{M/S} & \textbf{Policy} & \textbf{n} & \textbf{Cutoff} & \textbf{Worst recall} & \textbf{Mean recall} & \textbf{Worst AP} & \textbf{*} \\
\midrule\endfirsthead
\multicolumn{8}{l}{\textit{Table \thetable\ continued}}\\
\toprule
\textbf{M/S} & \textbf{Policy} & \textbf{n} & \textbf{Cutoff} & \textbf{Worst recall} & \textbf{Mean recall} & \textbf{Worst AP} & \textbf{*} \\
\midrule\endhead
\midrule\multicolumn{8}{r}{\textit{Continued on next page}}\\\endfoot
\bottomrule\endlastfoot
M & scene linear & 1 & 0.60429 & 0.0\% & 50.6\% & 0.3733 &  \\
M & region linear & 1 & 0.58476 & 0.0\% & 54.3\% & 0.4544 &  \\
M & region temporal seed 42 & 1 & 0.59187 & 5.3\% & 64.6\% & 0.5427 &  \\
M & region temporal seed 43 & 1 & 0.88623 & 12.7\% & 62.0\% & 0.5162 &  \\
M & region temporal seed 44 & 1 & 0.56992 & 1.5\% & 61.2\% & 0.5173 &  \\
M & region track temporal seed 42 & 1 & 0.40928 & 12.2\% & 61.2\% & 0.5492 &  \\
M & region track temporal seed 43 & 1 & 0.59090 & 12.7\% & 60.1\% & 0.5274 &  \\
M & region track temporal seed 44 & 1 & 0.64229 & 8.1\% & 64.1\% & 0.5175 &  \\
M & person supervised seed 42 & 1 & 0.68699 & 15.8\% & 56.6\% & 0.5704 &  \\
M & person supervised seed 43 & 1 & 0.71093 & 10.2\% & 56.5\% & 0.5723 &  \\
M & person supervised seed 44 & 1 & 0.69508 & 11.7\% & 57.2\% & 0.5723 &  \\
M & person supervised hardneg seed 42 & 1 & 0.67886 & 16.5\% & 58.7\% & 0.5736 & * \\
M & person supervised hardneg seed 43 & 1 & 0.68530 & 14.5\% & 56.3\% & 0.5670 &  \\
M & person supervised hardneg seed 44 & 1 & 0.67189 & 15.3\% & 56.2\% & 0.5719 &  \\
M & region temporal ensemble & 3 & 0.55311 & 11.5\% & 64.9\% & 0.5307 &  \\
M & region track temporal ensemble & 3 & 0.52465 & 12.5\% & 61.3\% & 0.5344 &  \\
M & person supervised ensemble & 3 & 0.69209 & 13.0\% & 57.8\% & 0.5711 &  \\
M & person supervised hardneg ensemble & 3 & 0.67862 & 15.5\% & 56.5\% & 0.5712 &  \\
M & all head ensemble & 14 & 0.52609 & 15.3\% & 65.8\% & 0.5717 &  \\
S & scene linear & 1 & 0.58306 & 0.0\% & 57.1\% & 0.3723 &  \\
S & region linear & 1 & 0.53811 & 0.0\% & 59.6\% & 0.4293 &  \\
S & region temporal seed 42 & 1 & 0.87444 & 15.0\% & 64.7\% & 0.5810 &  \\
S & region temporal seed 43 & 1 & 0.54643 & 13.5\% & 65.6\% & 0.5393 &  \\
S & region temporal seed 44 & 1 & 0.63610 & 13.7\% & 66.3\% & 0.5435 &  \\
S & region track temporal seed 42 & 1 & 0.75900 & 14.0\% & 65.7\% & 0.5710 &  \\
S & region track temporal seed 43 & 1 & 0.79908 & 13.7\% & 64.4\% & 0.5530 &  \\
S & region track temporal seed 44 & 1 & 0.85386 & 14.5\% & 62.1\% & 0.5774 &  \\
S & person supervised seed 42 & 1 & 0.57632 & 0.0\% & 69.2\% & 0.4921 &  \\
S & person supervised seed 43 & 1 & 0.63885 & 4.3\% & 70.3\% & 0.5531 &  \\
S & person supervised seed 44 & 1 & 0.62491 & 5.9\% & 70.7\% & 0.5542 &  \\
S & person supervised hardneg seed 42 & 1 & 0.53053 & 0.0\% & 68.9\% & 0.4890 &  \\
S & person supervised hardneg seed 43 & 1 & 0.63121 & 5.6\% & 70.6\% & 0.5471 &  \\
S & person supervised hardneg seed 44 & 1 & 0.61131 & 6.1\% & 70.8\% & 0.5506 &  \\
S & region temporal ensemble & 3 & 0.65598 & 14.2\% & 64.5\% & 0.5584 &  \\
S & region track temporal ensemble & 3 & 0.78296 & 15.8\% & 63.7\% & 0.5715 &  \\
S & person supervised ensemble & 3 & 0.61597 & 3.1\% & 70.0\% & 0.5497 &  \\
S & person supervised hardneg ensemble & 3 & 0.59184 & 4.8\% & 70.4\% & 0.5479 &  \\
S & all head ensemble & 14 & 0.54582 & 16.8\% & 66.4\% & 0.5789 & * \\
\end{longtable}\endgroup


\chapter{Detailed event results and source accounting}
This table retains global, historical, known-criteria and censored event outcomes from the pilot, including nonprimary candidates where the event policy was defined. Camera/scene counts and normal exposure are more informative than the number of overlapping decision windows. Zero normal exposure would make a normal hourly rate unavailable rather than zero. Unknown alerts are distinct from normal false alarms.
\begingroup\footnotesize
\setlength{\LTpre}{5pt}\setlength{\LTpost}{8pt}
\begin{longtable}{@{}>{\raggedright\arraybackslash}p{0.0928\linewidth}>{\raggedright\arraybackslash}p{0.2042\linewidth}>{\raggedright\arraybackslash}p{0.2320\linewidth}>{\raggedright\arraybackslash}p{0.0603\linewidth}>{\raggedright\arraybackslash}p{0.1021\linewidth}>{\raggedright\arraybackslash}p{0.0743\linewidth}>{\raggedright\arraybackslash}p{0.0743\linewidth}@{}}
\caption[Pilot all events]{Pilot event outcomes for all scored family routes; delay is median onset-to-emission time.}\label{tab:pilot_all_events}\\
\toprule
\textbf{Model} & \textbf{Protocol/cohort} & \textbf{Head} & \textbf{Hit / N} & \textbf{FP / dup / unk} & \textbf{Normal s} & \textbf{Delay s} \\
\midrule\endfirsthead
\multicolumn{7}{l}{\textit{Table \thetable\ continued}}\\
\toprule
\textbf{Model} & \textbf{Protocol/cohort} & \textbf{Head} & \textbf{Hit / N} & \textbf{FP / dup / unk} & \textbf{Normal s} & \textbf{Delay s} \\
\midrule\endhead
\midrule\multicolumn{7}{r}{\textit{Continued on next page}}\\\endfoot
\bottomrule\endlastfoot
Mobile\allowbreak{}CLIP & historical & scene linear & 1/5 & 2/0/3 & 230.0 & 0.0 \\
Mobile\allowbreak{}CLIP & fresh avss & scene linear & 3/3 & 6/6/0 & 264.1 & 23.2 \\
Mobile\allowbreak{}CLIP & fresh avss clean primary & scene linear & 2/2 & 4/5/0 & 165.7 & 20.6 \\
Mobile\allowbreak{}CLIP & fresh avss right censored diagnostic & scene linear & 1/1 & 2/1/0 & 98.3 & 23.2 \\
Mobile\allowbreak{}CLIP & fresh avss clean known 30 second rule & scene linear & 1/2 & 0/0/0 & 165.7 & 31.2 \\
Mobile\allowbreak{}CLIP & fresh avss clean state rule only & scene linear & 2/2 & 0/0/0 & 165.7 & 34.8 \\
Mobile\allowbreak{}CLIP & historical & region linear & 1/5 & 0/0/4 & 230.0 & 12.5 \\
Mobile\allowbreak{}CLIP & fresh avss & region linear & 3/3 & 4/4/0 & 264.1 & 18.2 \\
Mobile\allowbreak{}CLIP & fresh avss clean primary & region linear & 2/2 & 1/3/0 & 165.7 & 25.8 \\
Mobile\allowbreak{}CLIP & fresh avss right censored diagnostic & region linear & 1/1 & 3/1/0 & 98.3 & 18.2 \\
Mobile\allowbreak{}CLIP & fresh avss clean known 30 second rule & region linear & 2/2 & 0/0/0 & 165.7 & 36.3 \\
Mobile\allowbreak{}CLIP & fresh avss clean state rule only & region linear & 2/2 & 0/0/0 & 165.7 & 34.8 \\
Mobile\allowbreak{}CLIP & historical & region temporal seed 44 & 1/5 & 0/0/1 & 230.0 & 3.5 \\
Mobile\allowbreak{}CLIP & fresh avss & region temporal seed 44 & 3/3 & 1/1/0 & 264.1 & 55.2 \\
Mobile\allowbreak{}CLIP & fresh avss clean primary & region temporal seed 44 & 2/2 & 1/1/0 & 165.7 & 37.6 \\
Mobile\allowbreak{}CLIP & fresh avss right censored diagnostic & region temporal seed 44 & 1/1 & 0/0/0 & 98.3 & 88.2 \\
Mobile\allowbreak{}CLIP & fresh avss clean known 30 second rule & region temporal seed 44 & 1/2 & 0/0/0 & 165.7 & 44.9 \\
Mobile\allowbreak{}CLIP & fresh avss clean state rule only & region temporal seed 44 & 2/2 & 0/0/0 & 165.7 & 34.8 \\
Mobile\allowbreak{}CLIP & historical & region track temporal seed 44 & 1/5 & 0/0/1 & 230.0 & 7.0 \\
Mobile\allowbreak{}CLIP & fresh avss & region track temporal seed 44 & 0/3 & 0/0/0 & 264.1 & -- \\
Mobile\allowbreak{}CLIP & fresh avss clean primary & region track temporal seed 44 & 0/2 & 0/0/0 & 165.7 & -- \\
Mobile\allowbreak{}CLIP & fresh avss right censored diagnostic & region track temporal seed 44 & 0/1 & 0/0/0 & 98.3 & -- \\
Mobile\allowbreak{}CLIP & fresh avss clean known 30 second rule & region track temporal seed 44 & 0/2 & 0/0/0 & 165.7 & -- \\
Mobile\allowbreak{}CLIP & fresh avss clean state rule only & region track temporal seed 44 & 2/2 & 0/0/0 & 165.7 & 34.8 \\
Mobile\allowbreak{}CLIP & historical & grounded track temporal seed 43 & 1/5 & 5/0/5 & 230.0 & 7.5 \\
Mobile\allowbreak{}CLIP & fresh avss & grounded track temporal seed 43 & 2/3 & 0/0/0 & 264.1 & 37.1 \\
Mobile\allowbreak{}CLIP & fresh avss clean primary & grounded track temporal seed 43 & 2/2 & 0/0/0 & 165.7 & 37.1 \\
Mobile\allowbreak{}CLIP & fresh avss right censored diagnostic & grounded track temporal seed 43 & 0/1 & 0/0/0 & 98.3 & -- \\
Mobile\allowbreak{}CLIP & fresh avss clean known 30 second rule & grounded track temporal seed 43 & 0/2 & 0/0/0 & 165.7 & -- \\
Mobile\allowbreak{}CLIP & fresh avss clean state rule only & grounded track temporal seed 43 & 2/2 & 0/0/0 & 165.7 & 34.8 \\
SigLIP & historical & scene linear & 0/5 & 0/0/0 & 230.0 & -- \\
SigLIP & fresh avss & scene linear & 2/3 & 0/0/0 & 264.1 & 34.3 \\
SigLIP & fresh avss clean primary & scene linear & 1/2 & 0/0/0 & 165.7 & 18.4 \\
SigLIP & fresh avss right censored diagnostic & scene linear & 1/1 & 0/0/0 & 98.3 & 50.2 \\
SigLIP & fresh avss clean known 30 second rule & scene linear & 0/2 & 0/0/0 & 165.7 & -- \\
SigLIP & fresh avss clean state rule only & scene linear & 2/2 & 0/0/0 & 165.7 & 34.8 \\
SigLIP & historical & region linear & 1/5 & 1/0/5 & 230.0 & 1.0 \\
SigLIP & fresh avss & region linear & 2/3 & 4/1/0 & 264.1 & 62.8 \\
SigLIP & fresh avss clean primary & region linear & 1/2 & 3/1/0 & 165.7 & 47.9 \\
SigLIP & fresh avss right censored diagnostic & region linear & 1/1 & 1/0/0 & 98.3 & 77.7 \\
SigLIP & fresh avss clean known 30 second rule & region linear & 2/2 & 0/0/0 & 165.7 & 36.3 \\
SigLIP & fresh avss clean state rule only & region linear & 2/2 & 0/0/0 & 165.7 & 34.8 \\
SigLIP & historical & region temporal seed 42 & 0/5 & 5/0/9 & 230.0 & -- \\
SigLIP & fresh avss & region temporal seed 42 & 3/3 & 1/1/0 & 264.1 & 49.2 \\
SigLIP & fresh avss clean primary & region temporal seed 42 & 2/2 & 1/1/0 & 165.7 & 35.3 \\
SigLIP & fresh avss right censored diagnostic & region temporal seed 42 & 1/1 & 0/0/0 & 98.3 & 49.2 \\
SigLIP & fresh avss clean known 30 second rule & region temporal seed 42 & 0/2 & 0/0/0 & 165.7 & -- \\
SigLIP & fresh avss clean state rule only & region temporal seed 42 & 2/2 & 0/0/0 & 165.7 & 34.8 \\
SigLIP & historical & region track temporal seed 43 & 0/5 & 7/0/5 & 230.0 & -- \\
SigLIP & fresh avss & region track temporal seed 43 & 3/3 & 11/2/0 & 264.1 & 9.7 \\
SigLIP & fresh avss clean primary & region track temporal seed 43 & 2/2 & 7/1/0 & 165.7 & 5.6 \\
SigLIP & fresh avss right censored diagnostic & region track temporal seed 43 & 1/1 & 4/1/0 & 98.3 & 21.7 \\
SigLIP & fresh avss clean known 30 second rule & region track temporal seed 43 & 2/2 & 0/0/0 & 165.7 & 34.8 \\
SigLIP & fresh avss clean state rule only & region track temporal seed 43 & 2/2 & 0/0/0 & 165.7 & 34.8 \\
SigLIP & historical & grounded track temporal seed 42 & 0/5 & 6/0/3 & 230.0 & -- \\
SigLIP & fresh avss & grounded track temporal seed 42 & 3/3 & 9/5/0 & 264.1 & 17.9 \\
SigLIP & fresh avss clean primary & grounded track temporal seed 42 & 2/2 & 6/3/0 & 165.7 & 12.1 \\
SigLIP & fresh avss right censored diagnostic & grounded track temporal seed 42 & 1/1 & 3/2/0 & 98.3 & 21.7 \\
SigLIP & fresh avss clean known 30 second rule & grounded track temporal seed 42 & 2/2 & 0/1/0 & 165.7 & 34.8 \\
SigLIP & fresh avss clean state rule only & grounded track temporal seed 42 & 2/2 & 0/0/0 & 165.7 & 34.8 \\
\end{longtable}\endgroup

\begingroup\footnotesize
\setlength{\LTpre}{5pt}\setlength{\LTpost}{8pt}
\begin{longtable}{@{}>{\raggedright\arraybackslash}p{0.1500\linewidth}>{\raggedright\arraybackslash}p{0.2600\linewidth}>{\raggedright\arraybackslash}p{0.0600\linewidth}>{\raggedright\arraybackslash}p{0.1100\linewidth}>{\raggedright\arraybackslash}p{0.0800\linewidth}>{\raggedright\arraybackslash}p{0.0800\linewidth}>{\raggedright\arraybackslash}p{0.1000\linewidth}@{}}
\caption[Latest all event groups]{Latest state ablations by source/split; historical intervals have different source criteria.}\label{tab:latest_all_event_groups}\\
\toprule
\textbf{Source} & \textbf{Policy} & \textbf{Hit / N} & \textbf{FP / dup / unk} & \textbf{Normal s} & \textbf{Delay s} & \textbf{95\% upper/hr} \\
\midrule\endfirsthead
\multicolumn{7}{l}{\textit{Table \thetable\ continued}}\\
\toprule
\textbf{Source} & \textbf{Policy} & \textbf{Hit / N} & \textbf{FP / dup / unk} & \textbf{Normal s} & \textbf{Delay s} & \textbf{95\% upper/hr} \\
\midrule\endhead
\midrule\multicolumn{7}{r}{\textit{Continued on next page}}\\\endfoot
\bottomrule\endlastfoot
historical train & original tracks item memory & 0/4 & 1/0/2 & 78.00 & -- & 218.9 \\
historical train & robust tracks item memory & 0/4 & 1/0/2 & 78.00 & -- & 218.9 \\
historical train & original tracks global debounce & 0/4 & 1/0/1 & 78.00 & -- & 218.9 \\
historical train & robust tracks global debounce & 0/4 & 1/0/1 & 78.00 & -- & 218.9 \\
historical test & original tracks item memory & 4/5 & 0/1/1 & 230.00 & 19.24 & 46.9 \\
historical test & robust tracks item memory & 4/5 & 0/0/1 & 230.00 & 18.99 & 46.9 \\
historical test & original tracks global debounce & 4/5 & 0/0/1 & 230.00 & 19.24 & 46.9 \\
historical test & robust tracks global debounce & 4/5 & 0/0/1 & 230.00 & 18.99 & 46.9 \\
historical dev & original tracks item memory & 0/2 & 0/0/1 & 28.00 & -- & 385.2 \\
historical dev & robust tracks item memory & 0/2 & 0/0/1 & 28.00 & -- & 385.2 \\
historical dev & original tracks global debounce & 0/2 & 0/0/1 & 28.00 & -- & 385.2 \\
historical dev & robust tracks global debounce & 0/2 & 0/0/1 & 28.00 & -- & 385.2 \\
avss primary retrospective & original tracks item memory & 2/2 & 0/0/0 & 165.74 & 34.82 & 65.1 \\
avss primary retrospective & robust tracks item memory & 2/2 & 0/0/0 & 165.74 & 35.07 & 65.1 \\
avss primary retrospective & original tracks global debounce & 2/2 & 0/0/0 & 165.74 & 34.82 & 65.1 \\
avss primary retrospective & robust tracks global debounce & 2/2 & 0/0/0 & 165.74 & 35.07 & 65.1 \\
avss censored & original tracks item memory & 1/1 & 0/2/0 & 98.32 & 10.68 & 109.7 \\
avss censored & robust tracks item memory & 1/1 & 0/2/0 & 98.32 & 10.68 & 109.7 \\
avss censored & original tracks global debounce & 1/1 & 0/1/0 & 98.32 & 10.68 & 109.7 \\
avss censored & robust tracks global debounce & 1/1 & 0/1/0 & 98.32 & 10.68 & 109.7 \\
\end{longtable}\endgroup


\chapter{Complete SmolVLM histories and source results}
These tables retain all nine completed epoch records and all 60 source/method/seed summaries. O, J and R distinguish original, one joint adapter and per-task routing. Original rows repeat across seeds to support paired comparisons; they are the same frozen checkpoint. Rows and cohort counts are not independent events or extra cameras. Threshold definitions and retrospective qualifications are given in Chapter~\ref{ch:smol}. The bundled numeric CSV also retains every raw, development-F1 and development-budget field and TP/FP counts.
\begingroup\footnotesize
\setlength{\LTpre}{5pt}\setlength{\LTpost}{8pt}
\begin{longtable}{@{}>{\raggedright\arraybackslash}p{0.0869\linewidth}>{\raggedright\arraybackslash}p{0.1883\linewidth}>{\raggedright\arraybackslash}p{0.1014\linewidth}>{\raggedright\arraybackslash}p{0.1448\linewidth}>{\raggedright\arraybackslash}p{0.1883\linewidth}>{\raggedright\arraybackslash}p{0.1303\linewidth}@{}}
\caption[Smol history]{Complete epoch histories. Falling training loss does not imply better source transfer.}\label{tab:smol_history}\\
\toprule
\textbf{Seed} & \textbf{Learning rate} & \textbf{Ep.} & \textbf{Train loss} & \textbf{Dev macro AP} & \textbf{Visits} \\
\midrule\endfirsthead
\multicolumn{6}{l}{\textit{Table \thetable\ continued}}\\
\toprule
\textbf{Seed} & \textbf{Learning rate} & \textbf{Ep.} & \textbf{Train loss} & \textbf{Dev macro AP} & \textbf{Visits} \\
\midrule\endhead
\midrule\multicolumn{6}{r}{\textit{Continued on next page}}\\\endfoot
\bottomrule\endlastfoot
42 & 3e-05 & 1 & 0.5578 & 0.6584 & 7946 \\
42 & 3e-05 & 2 & 0.3964 & 0.6569 & 7946 \\
42 & 3e-05 & 3 & 0.3189 & 0.6565 & 7946 \\
42 & 0.0001 & 1 & 0.4789 & 0.7026 & 7946 \\
42 & 0.0001 & 2 & 0.2895 & 0.6485 & 7946 \\
42 & 0.0001 & 3 & 0.2235 & 0.6276 & 7946 \\
43 & 0.0001 & 1 & 0.4808 & 0.6548 & 7946 \\
43 & 0.0001 & 2 & 0.2878 & 0.6364 & 7946 \\
43 & 0.0001 & 3 & 0.2151 & 0.6216 & 7946 \\
\end{longtable}\endgroup

\begingroup\footnotesize
\setlength{\LTpre}{5pt}\setlength{\LTpost}{8pt}
\begin{longtable}{@{}>{\raggedright\arraybackslash}p{0.2348\linewidth}>{\raggedright\arraybackslash}p{0.0542\linewidth}>{\raggedright\arraybackslash}p{0.0542\linewidth}>{\raggedright\arraybackslash}p{0.0813\linewidth}>{\raggedright\arraybackslash}p{0.0813\linewidth}>{\raggedright\arraybackslash}p{0.0813\linewidth}>{\raggedright\arraybackslash}p{0.0903\linewidth}>{\raggedright\arraybackslash}p{0.0813\linewidth}>{\raggedright\arraybackslash}p{0.0813\linewidth}@{}}
\caption[Smol all metrics]{All 60 cohort/method/seed summaries. O=original, J=joint, R=per-task routed; AUC denotes AUROC. Rec./FPR use the dev-5\% point.}\label{tab:smol_all_metrics}\\
\toprule
\textbf{Source} & \textbf{Seed} & \textbf{M} & \textbf{AP} & \textbf{AUC} & \textbf{Raw F1} & \textbf{Dev-F1 F1} & \textbf{Rec.} & \textbf{FPR} \\
\midrule\endfirsthead
\multicolumn{9}{l}{\textit{Table \thetable\ continued}}\\
\toprule
\textbf{Source} & \textbf{Seed} & \textbf{M} & \textbf{AP} & \textbf{AUC} & \textbf{Raw F1} & \textbf{Dev-F1 F1} & \textbf{Rec.} & \textbf{FPR} \\
\midrule\endhead
\midrule\multicolumn{9}{r}{\textit{Continued on next page}}\\\endfoot
\bottomrule\endlastfoot
Smoking photos (Kaggle) & 42 & O & 0.9670 & 0.9619 & 0.6707 & 0.6667 & 93.8\% & 12.5\% \\
Smoking photos (Kaggle) & 42 & J & 0.9844 & 0.9754 & 0.9589 & 0.9636 & 93.8\% & 1.8\% \\
Smoking photos (Kaggle) & 42 & R & 0.9844 & 0.9754 & 0.9589 & 0.9636 & 93.8\% & 1.8\% \\
Smoking CCTV c1 & 42 & O & 0.5999 & 0.4835 & 0.1545 & 0.7214 & 0.0\% & 0.0\% \\
Smoking CCTV c1 & 42 & J & 0.6225 & 0.5016 & 0.0000 & 0.7537 & 1.8\% & 0.0\% \\
Smoking CCTV c1 & 42 & R & 0.6225 & 0.5016 & 0.0000 & 0.7537 & 1.8\% & 0.0\% \\
Littering / IWDD & 42 & O & 0.5922 & 0.5888 & 0.6614 & 0.6466 & 2.0\% & 4.0\% \\
Littering / IWDD & 42 & J & 0.7660 & 0.8032 & 0.6452 & 0.7895 & 30.0\% & 8.0\% \\
Littering / IWDD & 42 & R & 0.7660 & 0.8032 & 0.6452 & 0.7895 & 30.0\% & 8.0\% \\
Loitering / Street Scene & 42 & O & 0.4250 & 0.8452 & 0.0000 & 0.5161 & 0.0\% & 1.6\% \\
Loitering / Street Scene & 42 & J & 0.1822 & 0.4148 & 0.1333 & 0.1818 & 9.1\% & 6.2\% \\
Loitering / Street Scene & 42 & R & 0.1822 & 0.4148 & 0.1333 & 0.1818 & 9.1\% & 6.2\% \\
Public sleep / historical & 42 & O & 0.7437 & 0.8194 & 0.5000 & 0.6154 & 0.0\% & 0.0\% \\
Public sleep / historical & 42 & J & 0.6327 & 0.7917 & 0.6667 & 0.4000 & 33.3\% & 16.7\% \\
Public sleep / historical & 42 & R & 0.7885 & 0.8889 & 0.8000 & 0.4444 & 33.3\% & 8.3\% \\
Unattended / historical & 42 & O & 0.4762 & 0.6711 & 0.3721 & 0.4604 & 20.2\% & 1.9\% \\
Unattended / historical & 42 & J & 0.6933 & 0.8132 & 0.5368 & 0.4588 & 94.6\% & 70.8\% \\
Unattended / historical & 42 & R & 0.6933 & 0.8132 & 0.5368 & 0.4588 & 94.6\% & 70.8\% \\
Smoking / Stanford40 & 42 & O & 0.5831 & 0.8374 & 0.2991 & 0.2960 & 77.1\% & 25.4\% \\
Smoking / Stanford40 & 42 & J & 0.8900 & 0.9594 & 0.8118 & 0.7311 & 78.6\% & 4.4\% \\
Smoking / Stanford40 & 42 & R & 0.8900 & 0.9594 & 0.8118 & 0.7311 & 78.6\% & 4.4\% \\
Smoking / HMDB & 42 & O & 0.7638 & 0.9014 & 0.2943 & 0.2295 & 63.3\% & 2.9\% \\
Smoking / HMDB & 42 & J & 0.7837 & 0.9053 & 0.7487 & 0.6332 & 69.7\% & 2.8\% \\
Smoking / HMDB & 42 & R & 0.7837 & 0.9053 & 0.7487 & 0.6332 & 69.7\% & 2.8\% \\
Public sleep / new authors & 42 & O & 0.8062 & 0.7005 & 0.7556 & 0.7143 & 0.0\% & 0.0\% \\
Public sleep / new authors & 42 & J & 0.8441 & 0.7273 & 0.6429 & 0.5600 & 41.2\% & 0.0\% \\
Public sleep / new authors & 42 & R & 0.8904 & 0.7701 & 0.7742 & 0.5833 & 41.2\% & 0.0\% \\
MEVA abandonment action & 42 & O & 0.6076 & 0.5278 & 0.0000 & 0.4800 & 0.0\% & 0.0\% \\
MEVA abandonment action & 42 & J & 0.6457 & 0.5611 & 0.1176 & 0.7368 & 66.7\% & 58.3\% \\
MEVA abandonment action & 42 & R & 0.6457 & 0.5611 & 0.1176 & 0.7368 & 66.7\% & 58.3\% \\
Smoking photos (Kaggle) & 43 & O & 0.9670 & 0.9619 & 0.6707 & 0.6667 & 93.8\% & 12.5\% \\
Smoking photos (Kaggle) & 43 & J & 0.9848 & 0.9762 & 0.9636 & 0.9511 & 94.6\% & 1.8\% \\
Smoking photos (Kaggle) & 43 & R & 0.9848 & 0.9762 & 0.9636 & 0.9511 & 94.6\% & 1.8\% \\
Smoking CCTV c1 & 43 & O & 0.5999 & 0.4835 & 0.1545 & 0.7214 & 0.0\% & 0.0\% \\
Smoking CCTV c1 & 43 & J & 0.6295 & 0.5137 & 0.0000 & 0.7590 & 0.8\% & 0.0\% \\
Smoking CCTV c1 & 43 & R & 0.6295 & 0.5137 & 0.0000 & 0.7590 & 0.8\% & 0.0\% \\
Littering / IWDD & 43 & O & 0.5922 & 0.5888 & 0.6614 & 0.6466 & 2.0\% & 4.0\% \\
Littering / IWDD & 43 & J & 0.7656 & 0.8028 & 0.7748 & 0.7308 & 8.0\% & 0.0\% \\
Littering / IWDD & 43 & R & 0.7656 & 0.8028 & 0.7748 & 0.7308 & 8.0\% & 0.0\% \\
Loitering / Street Scene & 43 & O & 0.4250 & 0.8452 & 0.0000 & 0.5161 & 0.0\% & 1.6\% \\
Loitering / Street Scene & 43 & J & 0.1314 & 0.2983 & 0.0000 & 0.1017 & 0.0\% & 1.6\% \\
Loitering / Street Scene & 43 & R & 0.4250 & 0.8452 & 0.0000 & 0.5161 & 0.0\% & 1.6\% \\
Public sleep / historical & 43 & O & 0.7437 & 0.8194 & 0.5000 & 0.6154 & 0.0\% & 0.0\% \\
Public sleep / historical & 43 & J & 0.9167 & 0.9444 & 0.7692 & 0.5000 & 33.3\% & 0.0\% \\
Public sleep / historical & 43 & R & 0.9167 & 0.9444 & 0.7692 & 0.5000 & 33.3\% & 0.0\% \\
Unattended / historical & 43 & O & 0.4762 & 0.6711 & 0.3721 & 0.4604 & 20.2\% & 1.9\% \\
Unattended / historical & 43 & J & 0.7052 & 0.8048 & 0.6045 & 0.4112 & 94.2\% & 61.7\% \\
Unattended / historical & 43 & R & 0.7052 & 0.8048 & 0.6045 & 0.4112 & 94.2\% & 61.7\% \\
Smoking / Stanford40 & 43 & O & 0.5831 & 0.8374 & 0.2991 & 0.2960 & 77.1\% & 25.4\% \\
Smoking / Stanford40 & 43 & J & 0.8934 & 0.9619 & 0.7889 & 0.6564 & 81.4\% & 6.2\% \\
Smoking / Stanford40 & 43 & R & 0.8934 & 0.9619 & 0.7889 & 0.6564 & 81.4\% & 6.2\% \\
Smoking / HMDB & 43 & O & 0.7638 & 0.9014 & 0.2943 & 0.2295 & 63.3\% & 2.9\% \\
Smoking / HMDB & 43 & J & 0.7862 & 0.9028 & 0.7379 & 0.6585 & 69.7\% & 2.9\% \\
Smoking / HMDB & 43 & R & 0.7862 & 0.9028 & 0.7379 & 0.6585 & 69.7\% & 2.9\% \\
Public sleep / new authors & 43 & O & 0.8062 & 0.7005 & 0.7556 & 0.7143 & 0.0\% & 0.0\% \\
Public sleep / new authors & 43 & J & 0.8498 & 0.7433 & 0.6154 & 0.3810 & 23.5\% & 0.0\% \\
Public sleep / new authors & 43 & R & 0.8498 & 0.7433 & 0.6154 & 0.3810 & 23.5\% & 0.0\% \\
MEVA abandonment action & 43 & O & 0.6076 & 0.5278 & 0.0000 & 0.4800 & 0.0\% & 0.0\% \\
MEVA abandonment action & 43 & J & 0.6620 & 0.6056 & 0.1250 & 0.7317 & 53.3\% & 50.0\% \\
MEVA abandonment action & 43 & R & 0.6620 & 0.6056 & 0.1250 & 0.7317 & 53.3\% & 50.0\% \\
\end{longtable}\endgroup


\chapter{Pinned identities, artifact registry and acquisition references}
Full SHA-256 values are stored in the accompanying JSON manifest; the table shows compact prefixes for reading. Paths are repository-relative. The registry identifies what this documentation build read, including code inspected for implementation details. It is not a claim that every current code file is byte-identical to every original execution snapshot; those execution-specific identities remain in their original artifacts.
\begingroup\footnotesize
\setlength{\LTpre}{5pt}\setlength{\LTpost}{8pt}
\begin{longtable}{@{}>{\raggedright\arraybackslash}p{0.1273\linewidth}>{\raggedright\arraybackslash}p{0.3309\linewidth}>{\raggedright\arraybackslash}p{0.3818\linewidth}@{}}
\caption[Model revisions]{Pinned image-model identities from the adapter registry.}\label{tab:model_revisions}\\
\toprule
\textbf{Model} & \textbf{Repository} & \textbf{Revision} \\
\midrule\endfirsthead
\multicolumn{3}{l}{\textit{Table \thetable\ continued}}\\
\toprule
\textbf{Model} & \textbf{Repository} & \textbf{Revision} \\
\midrule\endhead
\midrule\multicolumn{3}{r}{\textit{Continued on next page}}\\\endfoot
\bottomrule\endlastfoot
Mobile\allowbreak{}CLIP & timm/MobileCLIP2-S2-OpenCLIP & \path{ac6b37c8fc40b62b623d09ed228a6f0c9bc29fe6} \\
SigLIP & google/siglip2-base-patch16-224 & \path{75de2d55ec2d0b4efc50b3e9ad70dba96a7b2fa2} \\
SmolVLM & HuggingFaceTB/SmolVLM-500M-Instruct & \path{a7da5b986cb59b408707209984f360a5f4ad7e47} \\
MiniCPM & openbmb/MiniCPM-V-4 & \path{6f69a8235885be89608f3a71ffd40695b379e1e7} \\
\end{longtable}\endgroup

\section{Preserved run-directory inventory}
These statuses are read from each root metadata file where present. ``Not recorded'' is not inferred to mean success. Prepared, interrupted, error and incomplete-data runs are preserved as diagnostics; the narrative identifies which completed studies provide the authoritative stage results. Separate source jobs assembled into a combined report are not extra independent datasets.
\begingroup\footnotesize
\setlength{\LTpre}{5pt}\setlength{\LTpost}{8pt}
\begin{longtable}{@{}>{\raggedright\arraybackslash}p{0.4504\linewidth}>{\raggedright\arraybackslash}p{0.1461\linewidth}>{\raggedright\arraybackslash}p{0.2435\linewidth}@{}}
\caption[Run inventory]{Preserved root run metadata status; not a combined performance leaderboard.}\label{tab:run_inventory}\\
\toprule
\textbf{Result directory} & \textbf{Metadata} & \textbf{Recorded status} \\
\midrule\endfirsthead
\multicolumn{3}{l}{\textit{Table \thetable\ continued}}\\
\toprule
\textbf{Result directory} & \textbf{Metadata} & \textbf{Recorded status} \\
\midrule\endhead
\midrule\multicolumn{3}{r}{\textit{Continued on next page}}\\\endfoot
\bottomrule\endlastfoot
\path{anomaly_headtohead} & run.json & not recorded \\
\path{clip_all_confirmed_20261013} & evaluation.json & interrupted \\
\path{clip_all_confirmed_20261013_v2} & evaluation.json & complete for available tasks \\
\path{clip_finetuning_20261015} & run.json & complete \\
\path{cpu_latency_20261013} & run.json & not recorded \\
\path{evaluation_20261011T104039_496754Z} & evaluation.json & interrupted \\
\path{evaluation_20261011T120031_096651Z} & evaluation.json & prepared \\
\path{evaluation_20261011T120249_353844Z} & evaluation.json & interrupted \\
\path{evaluation_20261011T120315_062384Z} & evaluation.json & complete for available tasks \\
\path{evaluation_20261012T034523_356317Z} & evaluation.json & prepared \\
\path{evaluation_20261012T035016_680032Z} & evaluation.json & interrupted \\
\path{evaluation_20261012T040547_431301Z} & evaluation.json & interrupted \\
\path{evaluation_20261012T040605_929193Z} & evaluation.json & error \\
\path{evaluation_20261012T165345_130231Z} & evaluation.json & error \\
\path{evaluation_20261013T061827_722009Z} & evaluation.json & error \\
\path{event_improvements_20261017} & run.json & complete \\
\path{expanded_accuracy_20261013} & state.json & complete \\
\path{five_task_readiness_20261013} & evaluation.json & incomplete data \\
\path{gpu_all_models_20261013} & evaluation.json & complete for available tasks \\
\path{gpu_full_accuracy_20261013} & evaluation.json & complete for available tasks \\
\path{gpu_minicpm_full_accuracy_20261013} & evaluation.json & complete for available tasks \\
\path{model_comparison_20261013} & evaluation.json & interrupted \\
\path{model_comparison_20261013_v2} & evaluation.json & complete for available tasks \\
\path{smolvlm_finetuning_20261019} & run.json & complete \\
\path{strategy_study_20261016} & run.json & complete \\
\path{video_pilot_20261016} & run.json & complete \\
\path{workflow_gpu_ready_pilot} & evaluation.json & complete for available tasks \\
\end{longtable}\endgroup

\section{Source registry}
\begingroup\footnotesize
\setlength{\LTpre}{5pt}\setlength{\LTpost}{8pt}
\begin{longtable}{@{}>{\raggedright\arraybackslash}p{0.0646\linewidth}>{\raggedright\arraybackslash}p{0.5874\linewidth}>{\raggedright\arraybackslash}p{0.1880\linewidth}@{}}
\caption[Source registry]{Documentation source registry; complete SHA-256 identities are in source\_manifest.json.}\label{tab:source_registry}\\
\toprule
\textbf{ID} & \textbf{Repository-relative source} & \textbf{SHA-256 prefix} \\
\midrule\endfirsthead
\multicolumn{3}{l}{\textit{Table \thetable\ continued}}\\
\toprule
\textbf{ID} & \textbf{Repository-relative source} & \textbf{SHA-256 prefix} \\
\midrule\endhead
\midrule\multicolumn{3}{r}{\textit{Continued on next page}}\\\endfoot
\bottomrule\endlastfoot
S001 & \path{anomaly_eval/config/anomalies.json} & 589882c0a43d846b \\
S002 & \path{anomaly_eval/models.py} & 799f40199789e3ba \\
S003 & \path{anomaly_eval/training/clip.py} & c400ce680f762926 \\
S004 & \path{anomaly_eval/training/data.py} & 5424b130aa53df77 \\
S005 & \path{anomaly_eval/training/event_memory.py} & 6b5150df1b6096cb \\
S006 & \path{anomaly_eval/training/smolvlm.py} & 7902175ed111964e \\
S007 & \path{anomaly_eval/training/strategies.py} & a92a363a55a4e5be \\
S008 & \path{anomaly_eval/training/temporal.py} & e1d62cc77bb64681 \\
S009 & \path{anomaly_eval/training/video.py} & ebefa15bb1206191 \\
S010 & \path{data/clip_training_20261015/commons_source_audit.json} & aef85b7a0c7b2620 \\
S011 & \path{data/clip_training_20261015/dataset.json} & 2abb33d2d852eeec \\
S012 & \path{data/clip_training_20261015/iwdd_source_audit.json} & e5b8a85233b1e3d1 \\
S013 & \path{data/clip_training_20261015/stanford_source_audit.json} & a32d853aeb257eab \\
S014 & \path{data/event_improvements_20261017/component_protocol.json} & 0547d4d78376b4bb \\
S015 & \path{data/event_improvements_20261017/receipt.json} & 3ea715d29628785f \\
S016 & \path{data/strategy_sources_20261016/meva/index.json} & 8825968c7e43a9d3 \\
S017 & \path{data/strategy_sources_20261016/small_sources_receipts.json} & 693ad195b4b01efa \\
S018 & \path{data/strategy_study_20261016/dataset.json} & 738025df2f939460 \\
S019 & \path{data/video_pilot_20261016/data_verification.json} & 706494b7b35c53e6 \\
S020 & \path{data/video_pilot_20261016/detector.json} & 6901655eb0cdbebd \\
S021 & \path{data/video_pilot_20261016/fresh_cohort.json} & c759448c14a2444b \\
S022 & \path{data/video_pilot_20261016/fresh_recordings.jsonl} & 58b640e8e906b2d5 \\
S023 & \path{data/video_pilot_20261016/grounding_all.json} & 5d16644a3da50282 \\
S024 & \path{data/video_pilot_20261016/source_guard.json} & 371c7477a3c9ec0d \\
S025 & \path{docs/evaluation.md} & 00db02e3c7cf204e \\
S026 & \path{docs/event_improvements.md} & 9f38600c38313640 \\
S027 & \path{docs/smolvlm_finetuning.md} & ceb323f2f2cba81f \\
S028 & \path{docs/strategy_study.md} & 567bf7f347dfb6d0 \\
S029 & \path{docs/training.md} & c4ecc6fa6e3f55d4 \\
S030 & \path{docs/video_pilot.md} & 19b9277e57b3ae39 \\
S031 & \path{reports/anomaly_benchmark_report.md} & 55e7710474e7354f \\
S032 & \path{reports/archive/2026-10-11/evaluation_audit.md} & f077b00644d3eca8 \\
S033 & \path{reports/archive/2026-10-11/evaluation_validation_summary.json} & 4c7a8967184ffdc2 \\
S034 & \path{reports/clip_finetuning_findings_20261016.md} & 734201d3930c873a \\
S035 & \path{reports/clip_training_protocol_20261015.md} & 673c4c5b69f019f3 \\
S036 & \path{reports/data_readiness_20261013.json} & 1d2a3cf7ac3edf5c \\
S037 & \path{reports/event_improvements_findings_20261017.md} & f57c2e781e054a78 \\
S038 & \path{reports/hardware_20261013.json} & f2ffc6760ed3748b \\
S039 & \path{reports/slides/progress_validation_20261017.json} & 0edfb0dbb6ab41e1 \\
S040 & \path{reports/smolvlm_finetuning_findings_20261019.md} & ee68ae7119060393 \\
S041 & \path{reports/strategy_findings_20261016.md} & 940ab943c63d4435 \\
S042 & \path{reports/strategy_protocol_20261016.md} & 08261f1b12abda00 \\
S043 & \path{reports/vadclip_source_validation.json} & 6e1e8d99bdb8502b \\
S044 & \path{reports/video_pilot_findings_20261016.md} & 08d109020bcd93c8 \\
S045 & \path{reports/video_pilot_protocol_20261016.md} & 029ec05c5583a4a8 \\
S046 & \path{reports/workflow_review_20261013.md} & afe2cfb6ad54abf1 \\
S047 & \path{results/anomaly_headtohead/run.json} & 3a1942b128f568a1 \\
S048 & \path{results/clip_all_confirmed_20261013/evaluation.json} & c38190048b077759 \\
S049 & \path{results/clip_all_confirmed_20261013_v2/evaluation.json} & 2be4dfc8650c046c \\
S050 & \path{results/clip_finetuning_20261015/comparison.csv} & 27af4c8fb13878e2 \\
S051 & \path{results/clip_finetuning_20261015/run.json} & ac4f15585af8dd99 \\
S052 & \path{results/clip_finetuning_20261015/seed_42/mobileclip2-s2/float32_evaluation/results.json} & db0c73c60eb5c31a \\
S053 & \path{results/clip_finetuning_20261015/seed_42/mobileclip2-s2/selection.json} & 3e15fc97805d3488 \\
S054 & \path{results/clip_finetuning_20261015/seed_42/siglip2-b16-224/float32_evaluation/results.json} & 460b098feb1d3c02 \\
S055 & \path{results/clip_finetuning_20261015/seed_42/siglip2-b16-224/selection.json} & 62820ea45cee7d26 \\
S056 & \path{results/clip_finetuning_20261015/seed_43/mobileclip2-s2/float32_evaluation/results.json} & 3f60d6c625b27794 \\
S057 & \path{results/clip_finetuning_20261015/seed_43/mobileclip2-s2/selection.json} & b2594aceb8747804 \\
S058 & \path{results/clip_finetuning_20261015/seed_43/siglip2-b16-224/float32_evaluation/results.json} & 9888428d4c30e4a6 \\
S059 & \path{results/clip_finetuning_20261015/seed_43/siglip2-b16-224/selection.json} & a8507c66816b5fed \\
S060 & \path{results/clip_finetuning_20261015/seed_44/mobileclip2-s2/float32_evaluation/results.json} & afada61f1de46d1d \\
S061 & \path{results/clip_finetuning_20261015/seed_44/mobileclip2-s2/selection.json} & b07d6d538c932634 \\
S062 & \path{results/clip_finetuning_20261015/seed_44/siglip2-b16-224/float32_evaluation/results.json} & 59761c9f8cbc2b1d \\
S063 & \path{results/clip_finetuning_20261015/seed_44/siglip2-b16-224/selection.json} & 04a7c8f977522543 \\
S064 & \path{results/clip_finetuning_20261015/verification.json} & 510410415c81480e \\
S065 & \path{results/cpu_latency_20261013/minicpm-v4/summary.json} & b44cc7a7c41a9bea \\
S066 & \path{results/cpu_latency_20261013/mobileclip2-s2/summary.json} & cbe6858112ab9321 \\
S067 & \path{results/cpu_latency_20261013/run.json} & b9497cc0e967c09f \\
S068 & \path{results/cpu_latency_20261013/siglip2-b16-224/summary.json} & 8c20f68fccdccad3 \\
S069 & \path{results/cpu_latency_20261013/smolvlm-500m/summary.json} & 0419b64d3acf3026 \\
S070 & \path{results/evaluation_20261011T104039_496754Z/evaluation.json} & 8ea7b7f59e4de28b \\
S071 & \path{results/evaluation_20261011T120031_096651Z/evaluation.json} & 21385e876c00fd08 \\
S072 & \path{results/evaluation_20261011T120249_353844Z/evaluation.json} & c58cfd292ca83f0c \\
S073 & \path{results/evaluation_20261011T120315_062384Z/evaluation.json} & 878d8af398f241d5 \\
S074 & \path{results/evaluation_20261012T034523_356317Z/evaluation.json} & 08a7bdeb13385f2d \\
S075 & \path{results/evaluation_20261012T035016_680032Z/evaluation.json} & a0421cb1f1e34ab2 \\
S076 & \path{results/evaluation_20261012T040547_431301Z/evaluation.json} & a0421cb1f1e34ab2 \\
S077 & \path{results/evaluation_20261012T040605_929193Z/evaluation.json} & 960841c767ff8966 \\
S078 & \path{results/evaluation_20261012T165345_130231Z/evaluation.json} & cc46d5af7ce79ae1 \\
S079 & \path{results/evaluation_20261012T165345_130231Z/report.md} & 4524008fae4559a5 \\
S080 & \path{results/evaluation_20261013T061827_722009Z/evaluation.json} & fcd5f73c8ae9714d \\
S081 & \path{results/event_improvements_20261017/component_correction.json} & 7214307a3ce826f2 \\
S082 & \path{results/event_improvements_20261017/component_results.json} & c17ed0fc199353d5 \\
S083 & \path{results/event_improvements_20261017/event_results.json} & 553119083ed95417 \\
S084 & \path{results/event_improvements_20261017/protocol.json} & ac309f57cacd6a4e \\
S085 & \path{results/event_improvements_20261017/report.md} & f57c2e781e054a78 \\
S086 & \path{results/event_improvements_20261017/run.json} & 2159e8ce72cd3855 \\
S087 & \path{results/event_improvements_20261017/smoking_results.json} & d363c41b3d186cc9 \\
S088 & \path{results/event_improvements_20261017/smoking_selection.json} & 9beab75476dab925 \\
S089 & \path{results/event_improvements_20261017/validation.json} & 7e1492e63b8da925 \\
S090 & \path{results/event_improvements_20261017/verification.json} & 8d47672cb508354c \\
S091 & \path{results/expanded_accuracy_20261013/quality_summary.json} & bd5adf87592a779b \\
S092 & \path{results/expanded_accuracy_20261013/state.json} & efaaa8bf6833b32a \\
S093 & \path{results/expanded_accuracy_20261013/verification.json} & ad313d179a83660f \\
S094 & \path{results/five_task_readiness_20261013/evaluation.json} & 7d77fa220dd048b0 \\
S095 & \path{results/gpu_all_models_20261013/evaluation.json} & 50b0b93b7ea0ed8c \\
S096 & \path{results/gpu_all_models_20261013/report.md} & 677e5c93d29144db \\
S097 & \path{results/gpu_all_models_20261013/smoking_frame/dev/minicpm-v4/summary.json} & 9aa833a2a3dc3dce \\
S098 & \path{results/gpu_all_models_20261013/smoking_frame/dev/mobileclip2-s2/summary.json} & 81ef7998a3c6ca16 \\
S099 & \path{results/gpu_all_models_20261013/smoking_frame/dev/siglip2-b16-224/summary.json} & 675e56ba96403504 \\
S100 & \path{results/gpu_all_models_20261013/smoking_frame/dev/smolvlm-500m/summary.json} & f0c9e86f86829442 \\
S101 & \path{results/gpu_all_models_20261013/unattended_items_frame/test/minicpm-v4/summary.json} & be005bd3154ea150 \\
S102 & \path{results/gpu_all_models_20261013/unattended_items_frame/test/mobileclip2-s2/summary.json} & 29f8f32074fbf431 \\
S103 & \path{results/gpu_all_models_20261013/unattended_items_frame/test/siglip2-b16-224/summary.json} & 24959f408523ed34 \\
S104 & \path{results/gpu_all_models_20261013/unattended_items_frame/test/smolvlm-500m/summary.json} & 1bd5994a794d33f0 \\
S105 & \path{results/gpu_all_models_20261013/unattended_items_window/test/minicpm-v4/summary.json} & a99de748065a6bca \\
S106 & \path{results/gpu_all_models_20261013/unattended_items_window/test/mobileclip2-s2/summary.json} & f1925a4e95608add \\
S107 & \path{results/gpu_all_models_20261013/unattended_items_window/test/siglip2-b16-224/summary.json} & 7f4793928d1041cc \\
S108 & \path{results/gpu_all_models_20261013/unattended_items_window/test/smolvlm-500m/summary.json} & bd0b8af8d51362a7 \\
S109 & \path{results/gpu_all_models_20261013/vadclip/adapter/eval.json} & 53480bc4b00d1d19 \\
S110 & \path{results/gpu_full_accuracy_20261013/evaluation.json} & 266956f97179412c \\
S111 & \path{results/gpu_minicpm_full_accuracy_20261013/evaluation.json} & c916084379862df9 \\
S112 & \path{results/model_comparison_20261013/evaluation.json} & 13c4e481f1354c75 \\
S113 & \path{results/model_comparison_20261013_v2/evaluation.json} & 8de485383b1b60b3 \\
S114 & \path{results/smolvlm_finetuning_20261019/additional_transfer/protocol.json} & 4a667a0a939383f6 \\
S115 & \path{results/smolvlm_finetuning_20261019/cli_checks.json} & b8833204438ad077 \\
S116 & \path{results/smolvlm_finetuning_20261019/clip_comparison_provenance.json} & 491cfd51323d7605 \\
S117 & \path{results/smolvlm_finetuning_20261019/environment.json} & 6f15970af209bfb0 \\
S118 & \path{results/smolvlm_finetuning_20261019/figures/learning_curves.pdf} & 837fdde5bb093791 \\
S119 & \path{results/smolvlm_finetuning_20261019/figures/matched_ap.pdf} & 8a50b3364f508a75 \\
S120 & \path{results/smolvlm_finetuning_20261019/implementation_manifest.json} & 482b8eb82dfaa060 \\
S121 & \path{results/smolvlm_finetuning_20261019/input_label_audit.json} & 21730eb54e96fe44 \\
S122 & \path{results/smolvlm_finetuning_20261019/metrics.csv} & e937cd9c7ec41ca1 \\
S123 & \path{results/smolvlm_finetuning_20261019/model_bundle.json} & c4f10d45feecda27 \\
S124 & \path{results/smolvlm_finetuning_20261019/portable_adapter_export.json} & d084a366d4c084ae \\
S125 & \path{results/smolvlm_finetuning_20261019/protocol.json} & e2b9ddc811daeceb \\
S126 & \path{results/smolvlm_finetuning_20261019/repeated_gain_protocol.json} & 7dfa0c31c058dc36 \\
S127 & \path{results/smolvlm_finetuning_20261019/repeated_gain_routes.json} & c8833fd6b23bd288 \\
S128 & \path{results/smolvlm_finetuning_20261019/report_data.json} & 91d32c25a1d71ee3 \\
S129 & \path{results/smolvlm_finetuning_20261019/response_mass.json} & add9c85dc8d0dcf3 \\
S130 & \path{results/smolvlm_finetuning_20261019/run.json} & b66a427586d11455 \\
S131 & \path{results/smolvlm_finetuning_20261019/seed_42/lr_0.0001/adapter/adapter_config.json} & 7c5e157b0b455b43 \\
S132 & \path{results/smolvlm_finetuning_20261019/seed_42/lr_0.0001/history.json} & 15b429993b498eeb \\
S133 & \path{results/smolvlm_finetuning_20261019/seed_42/lr_0.0001/state.json} & e90ab11f1a040465 \\
S134 & \path{results/smolvlm_finetuning_20261019/seed_42/lr_3e-05/adapter/adapter_config.json} & 7c5e157b0b455b43 \\
S135 & \path{results/smolvlm_finetuning_20261019/seed_42/lr_3e-05/history.json} & c8d46af1aad1ddd2 \\
S136 & \path{results/smolvlm_finetuning_20261019/seed_42/lr_3e-05/state.json} & bf1a56f43765bf43 \\
S137 & \path{results/smolvlm_finetuning_20261019/seed_42/selection.json} & a6025c69113aee4f \\
S138 & \path{results/smolvlm_finetuning_20261019/seed_43/lr_0.0001/adapter/adapter_config.json} & a66ce50439265a87 \\
S139 & \path{results/smolvlm_finetuning_20261019/seed_43/lr_0.0001/history.json} & f5e08567d711bc3e \\
S140 & \path{results/smolvlm_finetuning_20261019/seed_43/lr_0.0001/state.json} & 9b82c02690e0fdc3 \\
S141 & \path{results/smolvlm_finetuning_20261019/seed_43/selection.json} & 0259512a4409921e \\
S142 & \path{results/smolvlm_finetuning_20261019/software_checks.json} & 7b3e0c1688599464 \\
S143 & \path{results/smolvlm_finetuning_20261019/verification.json} & 17de97649f6789d6 \\
S144 & \path{results/strategy_study_20261016/comparison.csv} & 1451de20fec4a23d \\
S145 & \path{results/strategy_study_20261016/mobileclip2-s2/development_results.json} & 2ec38b5e140bcc88 \\
S146 & \path{results/strategy_study_20261016/mobileclip2-s2/selection.json} & 0cf75a9a614d3c6e \\
S147 & \path{results/strategy_study_20261016/mobileclip2-s2/transfer_evaluation/per_source_results.json} & 5462f52a7ecba313 \\
S148 & \path{results/strategy_study_20261016/report_data.json} & f2ad1868f968373a \\
S149 & \path{results/strategy_study_20261016/run.json} & 31f0f1424a033048 \\
S150 & \path{results/strategy_study_20261016/siglip2-b16-224/development_results.json} & 1aaa42385ea28801 \\
S151 & \path{results/strategy_study_20261016/siglip2-b16-224/selection.json} & 1aca5b02c95b3b41 \\
S152 & \path{results/strategy_study_20261016/siglip2-b16-224/transfer_evaluation/per_source_results.json} & fbfc44585ab6f8f3 \\
S153 & \path{results/strategy_study_20261016/thermal_trajectories/results.json} & fa3ce229a07b41fe \\
S154 & \path{results/strategy_study_20261016/verification.json} & ecd2984c5861159b \\
S155 & \path{results/vadclip/runs/full/eval.json} & b80eb22a5dfb82c9 \\
S156 & \path{results/vadclip/runs/no_graph/eval.json} & 6ce3846625a6d0b1 \\
S157 & \path{results/vadclip/runs/vl_128/eval.json} & d0e9386f5505e66a \\
S158 & \path{results/vadclip/runs/vl_512/eval.json} & 39dc97cc97f36ec1 \\
S159 & \path{results/vadclip/runs/win_1/eval.json} & beb67019aaabef2c \\
S160 & \path{results/vadclip/runs/win_16/eval.json} & 46feb1360d4d6ba9 \\
S161 & \path{results/video_pilot_20261016/comparison.csv} & c474b6adb4955ebf \\
S162 & \path{results/video_pilot_20261016/error_audit/pilot_region_coverage.json} & 57a619664b80ab33 \\
S163 & \path{results/video_pilot_20261016/events.csv} & cee168b1e111e7f3 \\
S164 & \path{results/video_pilot_20261016/mobileclip2-s2/event_policies.json} & 5f38904bccd14c9f \\
S165 & \path{results/video_pilot_20261016/mobileclip2-s2/joint_selection.json} & 3348807a3182eae0 \\
S166 & \path{results/video_pilot_20261016/run.json} & 9207ebb42c2803e0 \\
S167 & \path{results/video_pilot_20261016/siglip2-b16-224/joint_selection.json} & fa85f8a1cb9c59e6 \\
S168 & \path{results/video_pilot_20261016/verification.json} & 85a94ff384de4e75 \\
S169 & \path{results/workflow_gpu_ready_pilot/evaluation.json} & 6af247bed20d2a48 \\
S170 & \path{scripts/benchmark_item_tiles.py} & a5dd00d522e5c6b7 \\
S171 & \path{scripts/benchmark_loitering_trajectories.py} & 3295870fb8b00cf6 \\
S172 & \path{scripts/bootstrap_report_latex.py} & f7e23ec670d5325a \\
S173 & \path{scripts/build_technical_runthrough.py} & 22cf8e2de2ec9745 \\
S174 & \path{scripts/evaluate_video_pilot.py} & 377727fcd2ce65c0 \\
S175 & \path{scripts/fetch_item_component.py} & 3e45a01a59606a06 \\
S176 & \path{scripts/predict_smolvlm.py} & 52c88551955303ef \\
S177 & \path{scripts/prepare_item_component.py} & 8b06b0678db5b7d2 \\
S178 & \path{scripts/report_event_improvements.py} & 445563e442f08fd6 \\
S179 & \path{scripts/run_event_improvements.py} & e47654a9f4c27216 \\
S180 & \path{scripts/run_smolvlm_finetuning.py} & 75afb3ac09195f37 \\
S181 & \path{scripts/run_strategy_screen.py} & 3e64f3363f1ecda9 \\
S182 & \path{scripts/smolvlm_gpu.sh} & bf97aa57e9945dff \\
S183 & \path{scripts/train_grounded_pilot.py} & 03b88ff04f6d2e60 \\
S184 & \path{scripts/train_person_pilot.py} & f0fedd3b7731c799 \\
S185 & \path{scripts/train_video_pilot.py} & 52db372eed7ea451 \\
S186 & \path{scripts/verify_smolvlm_finetuning.py} & 83564369a4109fc1 \\
\end{longtable}\endgroup


\section{Recorded public source locations}
These URLs identify source releases, author pages or model repositories recorded by the project. They are not a new survey of current access conditions, and original dataset label/license metadata remains authoritative.
\begin{enumerate}
\item Stanford40 action images: \url{http://vision.stanford.edu/Datasets/40actions.html}.
\item MIVIA-IWDD author release: \url{https://github.com/MiviaLab/IllegalWasteDumping}.
\item Street Scene source: \url{https://www.merl.com/research/?research=license-request&sw=StreetScene}.
\item Added smoking photographs: \url{https://www.kaggle.com/datasets/sujaykapadnis/smoking}.
\item Recorded HMDB archive source: \url{https://huggingface.co/datasets/Serrelab/hmdb51}.
\item Thermal loitering annotations: \url{https://github.com/johnnynunez/RS-WACV24_Loitering}.
\item MEVA public data: \url{https://mevadata.org/}.
\item Commons media and per-image metadata: \url{https://commons.wikimedia.org/}.
\item CAVIAR source: \url{https://homepages.inf.ed.ac.uk/rbf/CAVIAR/}.
\item ABODA source: \url{https://github.com/kevinlin311tw/ABODA}.
\item AVSS/AOD author package: \url{http://www-vpu.eps.uam.es/publications/AODsurvey/}.
\item Additional author debris release: \url{https://github.com/hangsuuuu/video-of-abandoned-object-detection}.
\item Grounding DINO tiny weights: \url{https://huggingface.co/IDEA-Research/grounding-dino-tiny}.
\item Official Tectonic installation/release documentation: \url{https://tectonic-typesetting.github.io/en-US/install.html}.
\end{enumerate}

\section{Final scope statement}
The work produced standardized source-qualified evaluations, multiple trained research bundles, stronger representations/components and a more defensible account of generalization failures. It did not complete independent human verification of all target labels, an untouched multi-camera benchmark covering all five events, or a production streaming application. Those distinctions define what the existing results can support.
